\chapter{Reciprocal Avoidance: RVO and ORCA}
\label{ch:ch13}
% Function names for pseudocode are prefixed with "Orca" to stay unique
% across the book.
\SetKwFunction{OrcaHalfPlane}{OrcaHalfPlane}
\SetKwFunction{OrcaClosest}{OrcaClosestBoundaryPoint}
\SetKwFunction{OrcaLP}{OrcaLinearProgram}
\SetKwFunction{OrcaLPLine}{OrcaLineProgram}
\SetKwFunction{OrcaDense}{OrcaDenseFallback}
\SetKwFunction{OrcaStep}{OrcaStep}

\begin{objectives}
\begin{itemize}
  \item State precisely why two agents that both avoid each other with plain velocity obstacles oscillate, and why letting each agent take half of the responsibility removes the oscillation.
  \item Define the reciprocal velocity obstacle $\RVO_{A|B}$, show that it is the velocity obstacle with its apex moved to $\tfrac12(\vel_A + \vel_B)$, and explain why its guarantee does not extend to more than two agents.
  \item Derive the \orca half-plane step by step: the truncated velocity obstacle, the current relative velocity, the smallest change $\vect{u}$ that reaches its boundary with the outward normal $\vect{n}$, and the half-plane $\set{\vel : (\vel - (\vel_A + \tfrac12\vect{u}))\cdot\vect{n} \ge 0}$, including the full case analysis for $\vect{u}$.
  \item Solve the per-agent program, the velocity closest to the preferred one inside all half-planes and the speed disc, with the incremental two-dimensional linear program, and handle the dense case in which the half-planes have an empty intersection.
  \item Adapt the construction to non-cooperative obstacles (full responsibility) and to three dimensions (half-spaces), and state exactly what \orca guarantees and what it does not.
  \item Implement \orca in \python, reproduce the worked example, and compare no avoidance, the velocity obstacle (VO) rule and \orca in a simulation.
\end{itemize}
\end{objectives}

% ---------------------------------------------------------------------
\section{Why this matters for a drone swarm}
\label{sec:ch13-motivation}

Eight drones of a survey swarm fly nominal paths planned by conflict-based search (\cbs,
\cref{ch:ch09}). The paths are conflict-free on the grid, but the drones do
not fly the grid: they track the paths with a velocity controller, they are
late or early by a second or two, the wind pushes them sideways, and a
delivery drone that belongs to nobody in the swarm cuts across the survey
area. In the last three seconds before two drones would touch, a global
replan is too slow and a static plan is worthless. What is needed is a
\emph{local} rule that every drone can evaluate a dozen times per second from
what it sees around it: the positions, velocities and sizes of its
neighbours. The rule must return a velocity that is safe for a short horizon,
close to the velocity the nominal path asks for, and computable in
microseconds.

Velocity obstacles (VO, \cref{ch:ch12}) give such a rule for one agent among
moving obstacles that do not react: choose any velocity outside the cone of
velocities that lead to a collision. The rule breaks down when the obstacles
are other drones running the same rule. Each drone sees the other swerve,
concludes that the danger is gone, swerves back, and the pair performs a
\emph{reciprocal dance}: a zig-zag of velocity changes at every time step,
sometimes ending in a collision that neither drone intended. This chapter
repairs the rule in two steps. The reciprocal velocity obstacle
(RVO)~\cite{vandenberg2008rvo} lets each agent take half of the avoidance and
removes the oscillation for a pair. Optimal reciprocal collision avoidance
(\orca)~\cite{vandenberg2011orca} turns the cone into a half-plane of
permitted velocities per neighbour, so that the velocity choice among any
number of neighbours becomes a tiny linear program with a provable pairwise
guarantee. \orca is the local safety layer of the hybrid architecture of
\cref{ch:ch24} and the subject of Week~6 of the study plan (\cref{ch:appA}).
Its priority in the training plan is \priority{Essential}; RVO is
\priority{High}.

The chapter proceeds as follows. \Cref{sec:ch13-intuition} shows the dance
and the idea of sharing the responsibility. \Cref{sec:ch13-problem} fixes the
notation, recalls the velocity obstacle, states the oscillation precisely and
defines RVO. \Cref{sec:ch13-orca} derives the \orca half-plane with a figure
for every step, including the cases of a non-cooperative obstacle and of three
dimensions. \Cref{sec:ch13-lp} solves the per-agent program and its
infeasible variant. \Cref{sec:ch13-example} works through one encounter with
every number, and \cref{sec:ch13-simulation} runs eight agents.
\Cref{sec:ch13-properties} proves the guarantee and delimits it, and the
remaining sections cover variants, implementation, the drone system and
exercises.

% ---------------------------------------------------------------------
\section{The idea in plain words}
\label{sec:ch13-intuition}

Two people walk towards each other in a corridor. If both step to their left
at the same moment, they are still on a collision course, so both step back,
and the well-known corridor dance begins. The dance stops when both understand
that the other one is also going to move: each steps \emph{half} as far as it
would have to step alone, and keeps that step. That is the whole content of
reciprocity. \Cref{fig:ch13-idea} shows the same story for two agents that
choose velocities. With a plain velocity obstacle, A at time $t$ sees B on a
collision course and swerves by the full amount that takes A out of B's cone;
B does the same. One time step later, each sees the other's new velocity, and
in the new relative velocity the cone no longer contains the preferred
velocity: both return to it, and at $t + 2\dt$ they are head-on again. With
reciprocity, A swerves only halfway, trusting B to do the other half. The
relative velocity changes by the full amount, so the pair is safe, and one
step later nothing has changed: the halfway velocities are still the best
permitted ones, and the agents keep them.

\begin{figure}[tb]
  \centering
  \inputfigure{ch13/idea}
  \caption{The reciprocal dance. Top three rows: with plain velocity
  obstacles both agents swerve, both return, both swerve again, because each
  reacts to a neighbour who has already reacted. Fourth row: with reciprocity
  each agent takes half of the swerve, the relative velocity is the same as
  if one agent had taken all of it, and at the next step the situation is
  unchanged. Bottom: the trajectories, a zig-zag against one smooth swerve.}
  \label{fig:ch13-idea}
\end{figure}

RVO writes this as a set: A may choose $\vel$ only if the velocity
$2\vel - \vel_A$, which is twice as far from A's current velocity as $\vel$,
lies outside the ordinary velocity obstacle. \orca goes one step further. It
does not ask for a set of velocities outside a cone, which is not convex and
makes the choice among several neighbours a combinatorial problem. Instead it
computes, for the pair $(A, B)$, the smallest change $\vect{u}$ of the
relative velocity that reaches the boundary of the cone, gives half of it to
each agent, and permits A every velocity on the far side of the line through
$\vel_A + \tfrac12\vect{u}$ perpendicular to $\vect{u}$: a half-plane. B
receives the mirror half-plane. Whatever the two agents choose in their
half-planes, their relative velocity ends up outside the cone. With many
neighbours A intersects all its half-planes and picks the velocity closest to
the one it prefers, which is a two-dimensional linear program that is solved
in microseconds.

\begin{keyidea}
Never avoid alone what the other agent will also avoid. RVO lets each agent
take half of the velocity change; \orca computes that change $\vect{u}$ once
for the pair, gives $\tfrac12\vect{u}$ to each agent as a half-plane of
permitted velocities, and lets every agent choose, inside all its half-planes
and its speed limit, the velocity closest to its preferred one. If the
half-planes intersect and all agents use the same horizon, no pair collides
within that horizon.\index{ORCA!key idea}\index{reciprocity}
\end{keyidea}

% ---------------------------------------------------------------------
\section{Problem statement and notation}
\label{sec:ch13-problem}

\subsection{Agents, relative motion and velocity obstacles}
\label{sec:ch13-vo}

Agents are discs in the plane, as in \cref{ch:ch02}. Agent $A$ has centre
$\pos_A \in \R^2$, radius $r_A$, current velocity $\vel_A$, a
\textbf{preferred velocity}\index{preferred velocity} $\vel_A^{\mathrm{pref}}$
(the velocity it would fly if it were alone, for instance towards the next
waypoint of its nominal path at cruise speed) and a speed limit
$v_{\max}$. For a pair $(A, B)$ we write
\begin{equation}
  \pos_{\mathrm{rel}} = \pos_B - \pos_A, \qquad
  \vel_{\mathrm{rel}} = \vel_A - \vel_B, \qquad
  R = r_A + r_B
  \label{eq:ch13-relative}
\end{equation}
for the relative position, the relative velocity and the
\textbf{combined radius}\index{combined radius}, the same symbol $R$ that
\cref{ch:ch12} uses.
Seen from $A$, agent $B$ is a point at $\pos_{\mathrm{rel}}$ that moves with
velocity $-\vel_{\mathrm{rel}}$, and the two discs overlap at time $t \ge 0$
exactly when $\norm{\pos_{\mathrm{rel}} - t\,\vel_{\mathrm{rel}}} < R$,
that is, when $t\,\vel_{\mathrm{rel}}$ lies in the open disc
$\disc{\pos_{\mathrm{rel}}}{R}$ of the Minkowski sum
(\cref{ch:ch02}).\index{Minkowski sum} We use open discs throughout: a relative
velocity along which the discs merely touch counts as safe, and the safety
margin that every implementation adds to the radii
(\cref{sec:ch13-implementation}) covers the difference.

\Cref{ch:ch12} develops the velocity obstacle in detail; we recall what this
chapter needs, in the relative frame, which is where \orca lives.

\begin{definition}[Velocity obstacle, truncated velocity obstacle]
\label{def:ch13-vo}
For a pair $(A, B)$ with $\norm{\pos_{\mathrm{rel}}} > R$, the
\textbf{velocity obstacle}\index{velocity obstacle} of $A$ induced by $B$ is
the set of relative velocities that lead to a collision at some future time,
\begin{equation}
  \VO_{A|B} = \set{\vel \in \R^2 : \exists\, t > 0,\ t\,\vel \in \disc{\pos_{\mathrm{rel}}}{R}},
  \label{eq:ch13-vo}
\end{equation}
and the \textbf{truncated velocity obstacle}\index{velocity obstacle!truncated}
with horizon $\ttc > 0$ is the set of relative velocities that lead to a
collision within time $\ttc$,
\begin{equation}
  \VO^{\ttc}_{A|B} = \set{\vel \in \R^2 : \exists\, t \in (0, \ttc],\ t\,\vel \in \disc{\pos_{\mathrm{rel}}}{R}}
  = \bigcup_{t \in (0,\ttc]} \disc{\pos_{\mathrm{rel}}/t}{R/t}.
  \label{eq:ch13-vo-tau}
\end{equation}
In absolute terms, $\vel_A$ is dangerous for $A$ given $\vel_B$ if
$\vel_A - \vel_B \in \VO_{A|B}$; the set of such $\vel_A$ is the translate
$\vel_B + \VO_{A|B}$, which \cref{ch:ch12} writes as $\VO_{A|B}(\vel_B)$.
\end{definition}

$\VO_{A|B}$ is the open cone with apex at the origin whose two boundary rays,
the \emph{legs}, are tangent to the disc $\disc{\pos_{\mathrm{rel}}}{R}$; its
axis is the direction of $\pos_{\mathrm{rel}}$, at angle $\varphi$ from the
$x$-axis, and its half-angle is
\begin{equation}
  \theta = \arcsin\frac{R}{\norm{\pos_{\mathrm{rel}}}}.
  \label{eq:ch13-halfangle}
\end{equation}
The truncated set $\VO^{\ttc}_{A|B}$ is the union of the discs
$\disc{\pos_{\mathrm{rel}}/t}{R/t}$, which shrink and move towards the apex as
$t$ grows. All of them are tangent to the two legs, and the smallest one,
$\disc{\pos_{\mathrm{rel}}/\ttc}{R/\ttc}$, cuts the tip off the cone
(\cref{fig:ch13-orca-construction}(a)). The boundary of $\VO^{\ttc}_{A|B}$
therefore consists of three pieces: the arc of the smallest disc that faces
the apex, between the two points where the legs touch it, and the two legs
from those tangent points outwards. Membership is a ray--disc test
(\cref{ch:ch02}): $\vel \in \VO^{\ttc}_{A|B}$ exactly when the ray
$t \mapsto t\,\vel$ enters $\disc{\pos_{\mathrm{rel}}}{R}$ at a time at most
$\ttc$.\index{time to collision}

\subsection{The oscillation problem, stated precisely}
\label{sec:ch13-dance}

The velocity obstacle rule of \cref{ch:ch12} is: at every time step, choose the
velocity closest to $\vel^{\mathrm{pref}}$ that lies outside the velocity
obstacles induced by the \emph{current} velocities of the neighbours. When
the neighbours are also agents running this rule, the current velocities they
are reacting to are about to change, and the rule chases a moving target.
\Cref{fig:ch13-idea} shows the mechanism; the following proposition states
it in the simplest symmetric case, which is also the case that occurs most
often, because two drones that fly towards each other along a shared corridor
are in exactly this situation.

\begin{proposition}[The reciprocal dance of velocity obstacles]
\label{thm:ch13-dance}
Let two agents of equal radius fly towards each other with
$\vel_A^{\mathrm{pref}} = -\vel_B^{\mathrm{pref}} = (s, 0)$, far enough apart
that the cone half-angle $\theta$ does not change noticeably during two
time steps, and let both update their velocities synchronously with the
velocity obstacle rule (untruncated cones). Suppose that at step $k$ both
fly their preferred velocities and that both pass on the same side. Then at
step $k+1$ both deviate by $\delta = 2 s \sin\theta$ perpendicular to the
cone leg, at step $k+2$ both are back at their preferred velocities, and the
pattern repeats with period two.
\end{proposition}
\begin{proof}
At step $k$ the apex of $A$'s cone is at $\vel_B = (-s, 0)$, and
$\vel_A^{\mathrm{pref}} = (s, 0)$ lies on the axis of the cone at distance
$2s$ from the apex, hence at perpendicular distance $2s\sin\theta = \delta$
inside each leg. The closest velocity outside the cone is
$\vel_A^{\mathrm{pref}} + \delta\,\vect{n}$, where $\vect{n}$ is the outward
unit normal of the chosen leg; by the symmetry of the configuration $B$
chooses $\vel_B^{\mathrm{pref}} - \delta\,\vect{n}$ (passing on the same side
means opposite normals in the world frame). At step $k+1$ the apex of $A$'s
cone has moved to $\vel_B^{\mathrm{pref}} - \delta\,\vect{n}$, so the signed
distance of $\vel_A^{\mathrm{pref}}$ to the leg, measured along $\vect{n}$,
has grown from $-\delta$ to $0$: the preferred velocity lies on the boundary
of the cone and is permitted, and it is the closest permitted velocity to
itself. Both agents return to their preferred velocities, which restores the
situation of step $k$.
\end{proof}

The relative velocity of the pair alternates between the colliding
$(2s, 0)$ and the safe $(2s, 0) + 2\delta\,\vect{n}$, so the agents make
lateral progress only every other step, and their velocity changes by
$\delta$ at every step. In the simulation of \cref{fig:ch13-dance} (sampled
velocities, $\dt = 0.1$\,s) the lateral velocity of agent A under the
velocity obstacle rule jumps back and forth between about $+0.05$ and
$-0.18$\,m/s for the whole approach (range $[-0.18, +0.05]$\,m/s), and its
velocity change reverses direction 49~times in 100 steps. For a drone this is a stream of contradictory commands to the
attitude controller. Truncating the cones does not help: the two agents
still react to each other's stale velocities. What helps is to make the
rule anticipate the reaction of the other agent.\index{reciprocal dance}
\index{oscillation!velocity obstacles}

\subsection{Reciprocal velocity obstacles}
\label{sec:ch13-rvo}

\begin{definition}[Reciprocal velocity obstacle]
\label{def:ch13-rvo}
The \textbf{reciprocal velocity obstacle}\index{reciprocal velocity obstacle|textbf}
of agent $A$ induced by agent $B$ is
\begin{equation}
  \RVO_{A|B}(\vel_B, \vel_A) = \set{\vel \in \R^2 : 2\vel - \vel_A \in \vel_B + \VO_{A|B}}.
  \label{eq:ch13-rvo}
\end{equation}
Agent $A$ may choose $\vel$ only if $\vel \notin \RVO_{A|B}(\vel_B, \vel_A)$,
that is, only if $\vel$ is the average of its current velocity $\vel_A$ and
some velocity $\vel' = 2\vel - \vel_A$ that lies outside the velocity
obstacle: $\vel = \tfrac12(\vel_A + \vel')$.
\end{definition}

The definition says that $A$ goes only halfway from its current velocity to
a safe one, and relies on $B$ to go the other half. The set has a simple
shape (\cref{fig:ch13-rvo}).

\begin{proposition}[The apex moves to the midpoint]
\label{thm:ch13-rvo-apex}
$\RVO_{A|B}(\vel_B, \vel_A) = \tfrac12(\vel_A + \vel_B) + \VO_{A|B}$: the
reciprocal velocity obstacle is the velocity obstacle cone translated so
that its apex lies at the average of the two current velocities.
\end{proposition}
\begin{proof}
$2\vel - \vel_A \in \vel_B + \VO_{A|B}$ holds exactly when
$\vel \in \tfrac12(\vel_A + \vel_B) + \tfrac12\VO_{A|B}$, and
$\tfrac12\VO_{A|B} = \VO_{A|B}$ because a cone with apex at the origin is
closed under scaling by positive factors.
\end{proof}

\begin{figure}[tb]
  \centering
  \inputfigure{ch13/rvo}
  \caption{Velocity space of agent A in a head-on encounter,
  $\vel_A = (1, 0)$ and $\vel_B = (-1, 0)$, cone half-angle $20^\circ$. The
  velocity obstacle (orange) has its apex at $\vel_B$; the reciprocal
  velocity obstacle (blue, dashed) is the same cone with the apex at
  $\tfrac12(\vel_A + \vel_B)$. The closest velocity outside the RVO is the
  midpoint between $\vel_A$ and the closest velocity $\vel'$ outside the VO:
  A takes half of the swerve.}
  \label{fig:ch13-rvo}
\end{figure}

For the truncated cone the scaling argument changes the horizon:
$\tfrac12\VO^{\ttc}_{A|B} = \VO^{2\ttc}_{A|B}$, because halving the disc
$\disc{\pos_{\mathrm{rel}}/t}{R/t}$ gives $\disc{\pos_{\mathrm{rel}}/2t}{R/2t}$,
so the reciprocal version of a truncated cone with horizon $\ttc$ is a
truncated cone with horizon $2\ttc$ at the midpoint apex
(\cref{exr:ch13-rvo-truncated}). Now repeat the argument of
\cref{thm:ch13-dance} with the new rule.

\begin{proposition}[RVO removes the dance]
\label{thm:ch13-rvo-stable}
In the setting of \cref{thm:ch13-dance}, let both agents choose at every
step the velocity closest to their preferred one outside their reciprocal
velocity obstacle, passing on the same side. Then the velocities chosen at
step $k+1$ satisfy $\vel_A - \vel_B \notin \VO_{A|B}$, and they are chosen
again at every later step while the geometry is unchanged.
\end{proposition}
\begin{proof}
At step $k$ the apex of $A$'s reciprocal cone is at the midpoint
$\vect{m} = \tfrac12(\vel_A^{\mathrm{pref}} + \vel_B^{\mathrm{pref}}) =
\vect{0}$, so $\vel_A^{\mathrm{pref}}$ lies at distance $s$ from the apex
along the axis and at perpendicular distance $s\sin\theta = \delta/2$ inside
the leg. $A$ chooses $\vel_A' = \vel_A^{\mathrm{pref}} + \tfrac{\delta}{2}\vect{n}$
and, by symmetry, $B$ chooses $\vel_B' = \vel_B^{\mathrm{pref}} -
\tfrac{\delta}{2}\vect{n}$. The new relative velocity
$\vel_A' - \vel_B' = (2s, 0) + \delta\,\vect{n}$ lies on the boundary of
$\VO_{A|B}$, hence outside the open cone. At step $k+1$ the midpoint
$\tfrac12(\vel_A' + \vel_B') = \vect{m}$ is unchanged, so the reciprocal
cone is unchanged; $\vel_A^{\mathrm{pref}}$ is still inside by
$\delta/2$, and the closest permitted velocity is again $\vel_A'$. The same
holds for $B$, and by induction for every later step.
\end{proof}

Van den Berg, Lin and Manocha prove both statements in general, not only in
the symmetric case: if both agents choose velocities outside their
reciprocal velocity obstacles and pass on the same side, their new relative
velocity is outside the velocity obstacle, and if in addition both choose
the closest such velocity to their preferred one, the chosen velocities are
kept at the next step~\cite{vandenberg2008rvo}. In the simulation of
\cref{fig:ch13-dance}, RVO reverses the direction of its velocity change once
in 100 steps against 49 times for the velocity obstacle, and the sum of all
velocity changes of agent A drops from $8.32$ to $1.44$\,m/s.

\begin{figure}[tb]
  \centering
  \inputfigure{ch13/dance}
  \caption{Two agents swapping places head-on ($\pos_A = (-5, 0)$,
  $\pos_B = (5, 0.5)$, radii $0.5$, $\ttc = 5$\,s, $\dt = 0.1$\,s), simulated
  with \code{ch13\_orca.py}. Left: trajectories under sampled VO, sampled
  RVO and \orca. Right: the lateral velocity of A. The VO agent alternates
  between swerving and returning at every step (49 reversals), the RVO agent
  settles after two steps, the \orca agent flies one small constant lateral
  velocity. The dotted grey curve is RVO in an almost symmetric instance
  (offset $0.1$): the agents disagree about the side and switch it 16 times
  before they settle, the reciprocal dance of RVO.}
  \label{fig:ch13-dance}
\end{figure}

\paragraph{Why RVO is not enough.}
The guarantee is a two-agent statement, and it rests on two assumptions that
fail in a crowd. First, it assumes that $B$'s reaction mirrors $A$'s. With a
third agent $C$ present, $B$ chooses the closest velocity outside
$\RVO_{B|A} \cup \RVO_{B|C}$, and that velocity may lie on the boundary of
$\RVO_{B|C}$ alone: $B$ then takes none of the responsibility for $A$, the
relative velocity of $A$ and $B$ moves only halfway out of the cone, and the
pair is still on a collision course. Second, it assumes that both agents pass
on the same side. The union of several cones is not convex, and the closest
velocity outside it jumps from one gap to another when a cone moves a little;
two agents that choose different sides do not change their relative velocity
at all, they only move the midpoint, and at the next step one or both switch
sides. This is the \textbf{reciprocal dance}\index{reciprocal dance!of RVO}
of RVO, visible in the dotted curve of \cref{fig:ch13-dance} even for two
agents whose configuration is almost symmetric, and it is the reason why
van den Berg et al.\ state that RVO offers no guarantee for more than two
agents~\cite{vandenberg2011orca}. The hybrid reciprocal velocity obstacle
(HRVO) of Snape et al.\ combines one leg of the VO with one leg of the RVO
so that passing on the wrong side costs more, which reduces the dance but
does not remove it~\cite{snape2011hrvo}. \orca removes both problems at the
root: the permitted set per neighbour is a half-plane, which is convex, and
the share of the responsibility is fixed by construction rather than
assumed.

% ---------------------------------------------------------------------
\section{The \orca construction}
\label{sec:ch13-orca}

\orca replaces the question ``which velocities are safe for $A$?'' by the
question ``what is the smallest change of the \emph{relative} velocity that
makes the pair safe for the next $\ttc$ seconds, and how is it shared?''.
The answer is built in four steps, one per panel of
\cref{fig:ch13-orca-construction}. Throughout, $\ttc$ is the
\textbf{time horizon}\index{time horizon!ORCA}: \orca promises nothing beyond
it, and every agent recomputes its velocity long before it elapses.

\begin{figure}[tb]
  \centering
  \inputfigure{ch13/orca-construction}
  \caption{The four steps of the \orca construction on the worked example
  ($\pos_{\mathrm{rel}} = (4, 0.5)$, $R = 1$, $\ttc = 4$\,s). (a) The
  truncated velocity obstacle in relative-velocity space: a cone with apex
  at the origin, half-angle $\theta$ around the direction of
  $\pos_{\mathrm{rel}}$, cut off by the disc of centre $\pos_{\mathrm{rel}}/\ttc$
  and radius $R/\ttc$. (b) The current relative velocity
  $\vel_{\mathrm{rel}} = \vel_A - \vel_B$ lies inside. (c) $\vect{u}$ is the
  shortest vector from $\vel_{\mathrm{rel}}$ to the boundary, here to the
  right leg, and $\vect{n}$ the outward unit normal there. (d) In the
  velocity space of A, the permitted half-plane is bounded by the line through
  $\vel_A + \tfrac12\vect{u}$ perpendicular to $\vect{n}$; the hatched side
  is forbidden. The preferred velocity is forbidden and is projected onto the
  boundary. The dashed line is the boundary for full responsibility
  ($\vel_A + \vect{u}$), which applies against a non-cooperative obstacle.}
  \label{fig:ch13-orca-construction}
\end{figure}

\paragraph{Step 1: the truncated velocity obstacle.}
Form $\VO^{\ttc}_{A|B}$ of \cref{def:ch13-vo} in relative-velocity space
(\cref{fig:ch13-orca-construction}(a)). It depends only on
$\pos_{\mathrm{rel}}$, $R$ and $\ttc$, so both agents can compute the
same set: $B$ obtains $\VO^{\ttc}_{B|A} = -\VO^{\ttc}_{A|B}$, the mirror
image through the origin, because its relative position is
$-\pos_{\mathrm{rel}}$.

\paragraph{Step 2: the current relative velocity.}
Take $\vel_{\mathrm{rel}} = \vel_A - \vel_B$ from the current velocities
(\cref{fig:ch13-orca-construction}(b)). \orca calls the velocities around
which the construction is centred the \emph{optimisation velocities}; using
the current velocities is the choice that makes the half-planes of $A$ and
$B$ consistent, because both can observe them, and it is the choice of the
RVO2 library. If $\vel_{\mathrm{rel}} \in \VO^{\ttc}_{A|B}$, the agents collide
within $\ttc$ unless one of them changes velocity.

\paragraph{Step 3: the smallest change and the outward normal.}
Let $\vect{q}$ be the point of the boundary of $\VO^{\ttc}_{A|B}$ closest to
$\vel_{\mathrm{rel}}$, let
\begin{equation}
  \vect{u} = \vect{q} - \vel_{\mathrm{rel}}
  \label{eq:ch13-u}
\end{equation}
be the vector that carries $\vel_{\mathrm{rel}}$ to it, and let $\vect{n}$ be
the outward unit normal of $\VO^{\ttc}_{A|B}$ at $\vect{q}$
(\cref{fig:ch13-orca-construction}(c)). If $\vel_{\mathrm{rel}}$ is inside the
truncated cone, $\vect{u}$ points outward along $\vect{n}$ and
$\norm{\vect{u}}$ is the smallest velocity change of the pair that makes it
safe for $\ttc$. If $\vel_{\mathrm{rel}}$ is already outside, $\vect{u}$ points
towards the cone, against $\vect{n}$, and $\norm{\vect{u}}$ is the margin by
which the pair may drift towards danger. The boundary has three pieces, and
the closest point is found by the following case analysis, which is the
heart of every \orca implementation.

\begin{lemma}[Closest boundary point of the truncated velocity obstacle]
\label{thm:ch13-cases}
Let $\norm{\pos_{\mathrm{rel}}} > R$, let $\vect{c} = \pos_{\mathrm{rel}}/\ttc$
and $\rho = R/\ttc$ be the centre and radius of the truncating disc, and
let $\vect{w} = \vel_{\mathrm{rel}} - \vect{c}$. Write
$\pos \times \vect{w} = p_x w_y - p_y w_x$ for the two-dimensional cross
product and $\ell = \sqrt{\norm{\pos_{\mathrm{rel}}}^2 - R^2}$ for the length
of a tangent from the apex to $\disc{\pos_{\mathrm{rel}}}{R}$.
\begin{enumerate}[label=(\roman*)]
  \item \emph{Arc.} If $\vect{w}\cdot\pos_{\mathrm{rel}} < 0$ and
  $(\vect{w}\cdot\pos_{\mathrm{rel}})^2 > R^2\norm{\vect{w}}^2$, that is, if
  the angle between $\vect{w}$ and $-\pos_{\mathrm{rel}}$ is smaller than
  $\arccos(R/\norm{\pos_{\mathrm{rel}}}) = \pi/2 - \theta$, then the closest
  boundary point lies on the front arc:
  $\vect{n} = \vect{w}/\norm{\vect{w}}$ and
  $\vect{u} = (\rho - \norm{\vect{w}})\,\vect{n}$.
  \item \emph{Legs.} Otherwise the closest point is the foot of the
  perpendicular from $\vel_{\mathrm{rel}}$ onto one leg. The unit directions
  of the legs are $\vect{d}_{\mathrm{L}} = (p_x\ell - p_y R,\ p_x R + p_y\ell)/\norm{\pos_{\mathrm{rel}}}^2$
  (left) and $\vect{d}_{\mathrm{R}} = (p_x\ell + p_y R,\ -p_x R + p_y\ell)/\norm{\pos_{\mathrm{rel}}}^2$
  (right), that is, the direction of $\pos_{\mathrm{rel}}$ rotated by
  $+\theta$ and $-\theta$. If $\pos_{\mathrm{rel}} \times \vect{w} > 0$ the
  left leg is closest and $\vect{n} = (-d_{\mathrm{L},y}, d_{\mathrm{L},x})$;
  otherwise the right leg is closest and
  $\vect{n} = (d_{\mathrm{R},y}, -d_{\mathrm{R},x})$. In both cases
  $\vect{q} = (\vel_{\mathrm{rel}}\cdot\vect{d})\,\vect{d}$ with
  $\vect{d}$ the chosen leg and $\vect{u} = \vect{q} - \vel_{\mathrm{rel}}$.
\end{enumerate}
If $\norm{\pos_{\mathrm{rel}}} \le R$ the discs already overlap, the cone is
the whole plane, and the construction is applied instead to the disc
$\disc{\pos_{\mathrm{rel}}/\dt}{R/\dt}$ of relative velocities that leave the
discs overlapping after one time step: $\vect{w} = \vel_{\mathrm{rel}} -
\pos_{\mathrm{rel}}/\dt$, $\vect{n} = \vect{w}/\norm{\vect{w}}$ and
$\vect{u} = (R/\dt - \norm{\vect{w}})\,\vect{n}$.
\end{lemma}
\begin{proof}[Proof sketch]
The two legs are tangent to the truncating disc at the points where the
radii make the angle $\pi/2 - \theta$ with the direction from the centre
towards the apex, so the front arc spans exactly the directions $\vect{w}$
of case (i). For such $\vect{w}$ the nearest point of the circle is
$\vect{c} + \rho\,\vect{w}/\norm{\vect{w}}$, and no point of a leg is
nearer, because the legs lie outside the circle and touch it only at the
ends of the arc. For every other $\vect{w}$ the nearest boundary point is on
the leg on the same side of the axis as $\vect{w}$, and it is the foot of
the perpendicular, which lies beyond the tangent point precisely because
$\vect{w}$ is outside the front sector (\cref{exr:ch13-arc-case} asks for the
details). The overlap case is the ray--disc condition
$\norm{\pos_{\mathrm{rel}} - \dt\,\vel_{\mathrm{rel}}} \ge R$ rewritten as
$\vel_{\mathrm{rel}} \notin \disc{\pos_{\mathrm{rel}}/\dt}{R/\dt}$.
\end{proof}

The normal orientation deserves attention. In case (i), $\vect{n}$ points
away from the centre of the truncating disc, which is outward for the arc
whether $\vel_{\mathrm{rel}}$ is inside or outside the disc. In case (ii),
$\vect{n}$ is the leg direction rotated by $+90^\circ$ for the left leg and
by $-90^\circ$ for the right leg, which points out of the cone on either
side. The vector $\vect{u}$ is always parallel to $\vect{n}$; its sign says
whether $\vel_{\mathrm{rel}}$ is inside ($\vect{u}\cdot\vect{n} > 0$) or
outside ($\vect{u}\cdot\vect{n} < 0$) the truncated cone.\index{ORCA!normal orientation}

\paragraph{Step 4: the half-plane.}
Give half of the change to each agent (\cref{fig:ch13-orca-construction}(d)).

\begin{definition}[\orca half-plane]
\label{def:ch13-orca}
With $\vect{u}$ and $\vect{n}$ as in Step~3, the set of velocities permitted
to agent $A$ with respect to $B$ for the horizon $\ttc$ is the closed
half-plane\index{ORCA!half-plane|textbf}
\begin{equation}
  \ORCA^{\ttc}_{A|B} = \set{\vel \in \R^2 : \bigl(\vel - (\vel_A + \tfrac12\vect{u})\bigr)\cdot\vect{n} \ge 0}.
  \label{eq:ch13-orca}
\end{equation}
The half-plane of $B$ with respect to $A$ is obtained by the same
construction from $B$'s point of view; since $\VO^{\ttc}_{B|A} =
-\VO^{\ttc}_{A|B}$ and $\vel_B - \vel_A = -\vel_{\mathrm{rel}}$, its closest
point is $-\vect{q}$, its change is $-\vect{u}$ and its normal is
$-\vect{n}$, so
\begin{equation}
  \ORCA^{\ttc}_{B|A} = \set{\vel \in \R^2 : \bigl(\vel - (\vel_B - \tfrac12\vect{u})\bigr)\cdot(-\vect{n}) \ge 0}.
  \label{eq:ch13-orca-mirror}
\end{equation}
\end{definition}

The boundary of $\ORCA^{\ttc}_{A|B}$ passes through $\vel_A + \tfrac12\vect{u}$,
the current velocity shifted by half of the required change, and is
perpendicular to $\vect{u}$. Agent $A$ must move its velocity at least
$\tfrac12\norm{\vect{u}}$ in direction $\vect{n}$ when the pair is on a
collision course, and may move it at most $\tfrac12\norm{\vect{u}}$ against
$\vect{n}$ when it is not. This is why the method is called
\emph{reciprocal}: the responsibility is split half and half by construction,
and no assumption about the other agent's choice is needed beyond the fact
that it uses the same rule. It is called \emph{optimal} because the pair
$(\ORCA^{\ttc}_{A|B}, \ORCA^{\ttc}_{B|A})$ is what van den Berg et al.\ call
\emph{reciprocally maximal}\index{ORCA!reciprocal maximality}: among all pairs
of half-planes that are safe in the sense of \cref{thm:ch13-pairwise} and
treat the two agents symmetrically, no other pair permits more velocities
near the optimisation velocities; they argue this geometrically and note that
the equal split can be replaced by any fixed split that sums to
one~\cite{vandenberg2011orca}. \Cref{alg:ch13-halfplane}
collects Steps 1--4 with a general share $\alpha$.

\begin{algorithm}[htb]
\caption{The \orca half-plane of agent $A$ induced by neighbour $B$}
\label{alg:ch13-halfplane}
\KwIn{$\pos_A, \vel_A, r_A$ and $\pos_B, \vel_B, r_B$; horizon $\ttc$; time step $\dt$; responsibility share $\alpha$ ($\tfrac12$ reciprocal, $1$ non-cooperative)}
\KwOut{unit normal $\vect{n}$ and point $\vect{q}_A$ of the half-plane $\set{\vel : (\vel - \vect{q}_A)\cdot\vect{n} \ge 0}$}
\Fn{\OrcaHalfPlane{$A$, $B$, $\ttc$, $\dt$, $\alpha$}}{
  $\pos \gets \pos_B - \pos_A$; $\vel \gets \vel_A - \vel_B$; $R \gets r_A + r_B$\label{alg:ch13-halfplane:rel}\;
  $(\vect{u}, \vect{n}) \gets$ \OrcaClosest{$\pos$, $\vel$, $R$, $\ttc$, $\dt$}\;
  \KwRet{$(\vect{n},\ \vel_A + \alpha\,\vect{u})$}\label{alg:ch13-halfplane:ret}\;
}
\Fn{\OrcaClosest{$\pos$, $\vel$, $R$, $\ttc$, $\dt$}}{
  \If{$\norm{\pos} \le R$}{
    $\vect{c} \gets \pos/\dt$; $\rho \gets R/\dt$; $\vect{w} \gets \vel - \vect{c}$; $\vect{n} \gets \vect{w}/\norm{\vect{w}}$\label{alg:ch13-halfplane:overlap}\tcp*{overlapping: separate in one step}
    \KwRet{$((\rho - \norm{\vect{w}})\,\vect{n},\ \vect{n})$}\;
  }
  $\vect{c} \gets \pos/\ttc$; $\rho \gets R/\ttc$; $\vect{w} \gets \vel - \vect{c}$\label{alg:ch13-halfplane:disc}\tcp*{the disc that truncates the cone}
  \If{$\vect{w}\cdot\pos < 0$ \KwAnd $(\vect{w}\cdot\pos)^2 > R^2\,\norm{\vect{w}}^2$}{
    $\vect{n} \gets \vect{w}/\norm{\vect{w}}$\label{alg:ch13-halfplane:arc}\tcp*{front sector: closest point on the arc}
    \KwRet{$((\rho - \norm{\vect{w}})\,\vect{n},\ \vect{n})$}\;
  }
  $\ell \gets \sqrt{\norm{\pos}^2 - R^2}$\tcp*{length of a tangent from the apex}
  \eIf{$\pos \times \vect{w} > 0$}{
    $\vect{d} \gets \bigl(p_x\ell - p_y R,\ p_x R + p_y\ell\bigr)/\norm{\pos}^2$; $\vect{n} \gets (-d_y, d_x)$\label{alg:ch13-halfplane:left}\tcp*{left leg}
  }{
    $\vect{d} \gets \bigl(p_x\ell + p_y R,\ -p_x R + p_y\ell\bigr)/\norm{\pos}^2$; $\vect{n} \gets (d_y, -d_x)$\label{alg:ch13-halfplane:right}\tcp*{right leg}
  }
  $\vect{q} \gets (\vel\cdot\vect{d})\,\vect{d}$\tcp*{foot of the perpendicular on the leg}
  \KwRet{$(\vect{q} - \vel,\ \vect{n})$}\label{alg:ch13-halfplane:leg}\;
}
\end{algorithm}

\paragraph{Walkthrough.}
Line~\ref{alg:ch13-halfplane:rel} moves to the relative frame.
Line~\ref{alg:ch13-halfplane:overlap} is the emergency case of
\cref{thm:ch13-cases}: the discs overlap, the horizon is replaced by one time
step, and the half-plane pushes the agents apart. Line~\ref{alg:ch13-halfplane:disc}
forms the truncating disc and the vector $\vect{w}$ from its centre to the
relative velocity; the test that follows is case (i) of the lemma, with the
angle condition written without trigonometric functions. Lines
\ref{alg:ch13-halfplane:left}--\ref{alg:ch13-halfplane:right} select the leg
by the sign of the cross product and rotate its direction into the outward
normal; the formulas for $\vect{d}$ are the direction of $\pos$ rotated by
$\pm\theta$, using $\cos\theta = \ell/\norm{\pos}$ and
$\sin\theta = R/\norm{\pos}$. Line~\ref{alg:ch13-halfplane:leg} returns the
change to the foot of the perpendicular, and
line~\ref{alg:ch13-halfplane:ret} shifts the current velocity by the share
$\alpha$ of it.

\subsection{Non-cooperative obstacles: full responsibility}
\label{sec:ch13-noncooperative}

The half split is right only when $B$ runs the same rule. An intruder that
does not know about $A$, a static obstacle, or a drone whose controller has
failed will not take its half. If $A$ nevertheless takes only half of
$\vect{u}$, the relative velocity moves only halfway to the boundary and the
collision course remains. In that case $A$ takes \textbf{full
responsibility}\index{ORCA!full responsibility|textbf}\index{non-cooperative obstacle}:
\begin{equation}
  \ORCA^{\ttc,\,1}_{A|B} = \set{\vel \in \R^2 : \bigl(\vel - (\vel_A + \vect{u})\bigr)\cdot\vect{n} \ge 0},
  \label{eq:ch13-orca-full}
\end{equation}
the dashed line of \cref{fig:ch13-orca-construction}(d), obtained from
\cref{alg:ch13-halfplane} with $\alpha = 1$. For a static obstacle,
$\vel_B = \vect{0}$ and the relative velocity is $\vel_A$ itself; for an
intruder, $\vel_B$ is its estimated or predicted velocity
(\cref{ch:ch18,ch:ch20}), and $r_B$ is inflated by the uncertainty of the
prediction. In general, if $A$ and $B$ use shares $\alpha_A$ and $\alpha_B$,
the argument of \cref{thm:ch13-pairwise} below goes through exactly when
$\alpha_A + \alpha_B = 1$: the reciprocal split is $(\tfrac12, \tfrac12)$,
the non-cooperative case is $(1, 0)$, where $B$ keeps the velocity that $A$
assumed for it, and a priority scheme is any other split. This is the case
that matters for the hybrid architecture of \cref{ch:ch24}: the swarm mates
of a drone are reciprocal, the intruder is not, and the two kinds of
half-planes go into the same program.

\subsection{Three dimensions: half-spaces}
\label{sec:ch13-3d}

Drones fly in $\R^3$, and nothing in the construction was specific to the
plane. The velocity obstacle becomes the circular cone with apex at the
origin and half-angle $\theta = \arcsin(R/\norm{\pos_{\mathrm{rel}}})$ around
the axis $\pos_{\mathrm{rel}}$, truncated by the ball of centre
$\pos_{\mathrm{rel}}/\ttc$ and radius $R/\ttc$; its boundary is the front cap
of the sphere and the lateral surface of the cone. The truncated cone is a
body of revolution about the axis $\pos_{\mathrm{rel}}$, so the closest
boundary point to $\vel_{\mathrm{rel}}$ can be taken in the plane spanned by
$\pos_{\mathrm{rel}}$ and $\vel_{\mathrm{rel}}$ (the truncated cone is invariant
under reflection in any plane containing $\pos_{\mathrm{rel}}$, so the set of
nearest boundary points is invariant too; when $\vel_{\mathrm{rel}}$ lies
outside the body the nearest point of a convex set is unique and therefore
lies in the plane of $\pos_{\mathrm{rel}}$ and $\vel_{\mathrm{rel}}$, and when
it lies inside, at least one nearest boundary point lies in that plane and the
construction may take it; if $\vel_{\mathrm{rel}}$ is parallel to
$\pos_{\mathrm{rel}}$ every plane through the axis gives an equally good
$\vect{u}$ and the code picks one, exactly as the two-dimensional code breaks
the leg tie by a sign), and in that plane the body is exactly the
two-dimensional truncated cone. Hence $\vect{u}$ and $\vect{n}$ are computed by
\cref{thm:ch13-cases} in the coordinates of that plane and mapped back, and
the permitted set is the \textbf{half-space}\index{ORCA!three-dimensional}
\begin{equation}
  \ORCA^{\ttc}_{A|B} = \set{\vel \in \R^3 : \bigl(\vel - (\vel_A + \tfrac12\vect{u})\bigr)\cdot\vect{n} \ge 0}.
  \label{eq:ch13-orca-3d}
\end{equation}
The function \code{orca\_half\_space\_3d} of \code{ch13\_orca.py} does exactly
this, spanning the plane by the unit vector along $\pos_{\mathrm{rel}}$ and the
unit vector along the component of $\vel_{\mathrm{rel}}$ perpendicular to it,
or, when that component vanishes, along a fixed helper axis; its self-test
checks that the result is the two-dimensional
half-plane when everything lies in a plane and that it rotates with the
configuration. The per-agent program of the next section becomes a
three-dimensional linear program over half-spaces and the ball
$\norm{\vel} \le v_{\max}$; the incremental algorithm extends to three
dimensions by recursion (when a half-space is violated, solve the
two-dimensional program on its boundary plane) and keeps its expected linear
running time~\cite{seidel1991small,deberg2008computational}. The RVO2-3D
library implements this version.

% ---------------------------------------------------------------------
\section{The per-agent program}
\label{sec:ch13-lp}

\subsection{The velocity closest to the preferred one}

Agent $A$ with neighbours $B_1, \dots, B_m$ has $m$ half-planes
$H_i = \set{\vel : (\vel - \vect{q}_i)\cdot\vect{n}_i \ge 0}$ from
\cref{alg:ch13-halfplane}, and it must also respect its speed limit. It
chooses
\begin{equation}
  \vel_A^{\mathrm{new}} = \argmin_{\vel}\ \norm{\vel - \vel_A^{\mathrm{pref}}}
  \qquad \st \vel \in H_1 \cap \dots \cap H_m \cap \disc{\vect{0}}{v_{\max}}.
  \label{eq:ch13-program}
\end{equation}
The feasible set is an intersection of half-planes and a disc, hence
convex, and the objective is the distance to a point, so the program has a
unique solution whenever the feasible set is non-empty
(\cref{fig:ch13-lp}). It is a quadratic program, but a very special one: the
solution is either $\vel_A^{\mathrm{pref}}$ itself, or its projection onto
the disc, or it lies on the boundary line of one of the half-planes (or at
the intersection of two of them, or of one with the circle). This structure
is what the \textbf{incremental linear program}\index{linear program!incremental|textbf}
of \cref{alg:ch13-lp} exploits; it is the algorithm of
Seidel~\cite{seidel1991small} and of de Berg et al.~\cite{deberg2008computational},
specialised to two dimensions as in the RVO2 library.

\begin{figure}[tb]
  \centering
  \inputfigure{ch13/lp}
  \caption{The per-agent program in velocity space: the speed disc and three
  half-planes $H_1$, $H_2$, $H_3$ (hatched bands on the forbidden sides,
  green normals pointing into the permitted sides). The feasible set is
  convex. The preferred velocity violates $H_1$ only, and the solution is
  its projection onto the boundary line of $H_1$, which happens to satisfy
  $H_2$, $H_3$ and the disc. Had the projection violated $H_2$, the
  solution would sit at the intersection of the two boundary lines.}
  \label{fig:ch13-lp}
\end{figure}

\begin{algorithm}[htb]
\caption{Incremental two-dimensional linear program over half-planes and the speed disc}
\label{alg:ch13-lp}
\KwIn{half-planes $H_i = \set{\vel : (\vel - \vect{q}_i)\cdot\vect{n}_i \ge 0}$, $i = 1, \dots, m$; speed limit $v_{\max}$; preferred velocity $\vel^{\mathrm{pref}}$}
\KwOut{$(\vel, m+1)$ with $\vel$ the velocity closest to $\vel^{\mathrm{pref}}$ in $H_1 \cap \dots \cap H_m \cap \disc{\vect{0}}{v_{\max}}$, or $(\vel, i)$ if $H_i$ makes the intersection empty, $\vel$ then solving the program for $H_1, \dots, H_{i-1}$}
\Fn{\OrcaLP{$H_1, \dots, H_m$, $v_{\max}$, $\vel^{\mathrm{pref}}$}}{
  $\vel \gets \vel^{\mathrm{pref}}$, scaled to length $v_{\max}$ if $\norm{\vel^{\mathrm{pref}}} > v_{\max}$\label{alg:ch13-lp:init}\;
  \For{$i \gets 1$ \KwTo $m$}{
    \If{$(\vel - \vect{q}_i)\cdot\vect{n}_i < 0$}{\label{alg:ch13-lp:violated}
      $\vel' \gets$ \OrcaLPLine{$i$, $H_1, \dots, H_{i-1}$, $v_{\max}$, $\vel^{\mathrm{pref}}$}\tcp*{$H_i$ is violated: the optimum lies on its boundary}
      \lIf{$\vel' = $ \KwNil}{\KwRet{$(\vel, i)$}\label{alg:ch13-lp:fail}}
      $\vel \gets \vel'$\;
    }
  }
  \KwRet{$(\vel, m+1)$}\;
}
\Fn{\OrcaLPLine{$i$, $H_1, \dots, H_{i-1}$, $v_{\max}$, $\vel^{\mathrm{pref}}$}}{
  $\vect{d} \gets (n_{i,y}, -n_{i,x})$\tcp*{along the boundary line, with $H_i$ on its left}
  $b \gets \vect{q}_i\cdot\vect{d}$; $\Delta \gets b^2 + v_{\max}^2 - \norm{\vect{q}_i}^2$\;
  \lIf{$\Delta < 0$}{\KwRet{\KwNil} (the line misses the disc)}
  $t_- \gets -b - \sqrt{\Delta}$; $t_+ \gets -b + \sqrt{\Delta}$\label{alg:ch13-lp:chord}\tcp*{chord of the line inside the disc}
  \For{$j \gets 1$ \KwTo $i-1$}{
    $\vect{d}_j \gets (n_{j,y}, -n_{j,x})$; $\mathit{den} \gets \vect{d} \times \vect{d}_j$; $\mathit{num} \gets \vect{d}_j \times (\vect{q}_i - \vect{q}_j)$\;
    \If{$\abs{\mathit{den}} \le \eps$ \textnormal{($\eps = 10^{-9}$, a numerical tolerance)}}{\lIf{$\mathit{num} < 0$}{\KwRet{\KwNil} (parallel, $H_j$ excludes the whole line)} \lElse{\KwContinue}}
    \lIf{$\mathit{den} > 0$}{$t_+ \gets \min(t_+, \mathit{num}/\mathit{den})$} \lElse{$t_- \gets \max(t_-, \mathit{num}/\mathit{den})$\label{alg:ch13-lp:clip}}
    \lIf{$t_- > t_+$}{\KwRet{\KwNil}}
  }
  $t \gets \min\bigl(\max(\vect{d}\cdot(\vel^{\mathrm{pref}} - \vect{q}_i),\ t_-),\ t_+\bigr)$\label{alg:ch13-lp:clamp}\;
  \KwRet{$\vect{q}_i + t\,\vect{d}$}\;
}
\end{algorithm}

\paragraph{Walkthrough.}
The invariant of \cref{alg:ch13-lp} is that after the $i$-th iteration,
$\vel$ is the optimum of the program with the first $i$ half-planes. Line
\ref{alg:ch13-lp:init} solves the program with no half-planes: the preferred
velocity, clipped to the disc. If the next half-plane $H_i$ is satisfied by
the current optimum, adding it changes nothing, because the feasible set
only shrinks (line~\ref{alg:ch13-lp:violated}). If $H_i$ is violated, the
new optimum must lie on the boundary line of $H_i$: the old optimum was the
closest feasible point to $\vel^{\mathrm{pref}}$ in a convex set, the new
feasible set is the old one cut by $H_i$, and by convexity the closest point
of the cut set cannot lie strictly inside $H_i$ (\cref{exr:ch13-lp-by-hand}).
So the problem drops to one dimension: \OrcaLPLine parameterises the line as
$\vect{q}_i + t\,\vect{d}$, finds the interval $[t_-, t_+]$ of the line inside
the disc (line~\ref{alg:ch13-lp:chord}, from the quadratic
$\norm{\vect{q}_i + t\vect{d}}^2 = v_{\max}^2$), clips this interval by
every earlier half-plane (line~\ref{alg:ch13-lp:clip}: the point
$\vect{q}_i + t\vect{d}$ satisfies $H_j$ when $t\cdot\mathit{den} \le
\mathit{num}$), and takes the parameter closest to the projection of
$\vel^{\mathrm{pref}}$ (line~\ref{alg:ch13-lp:clamp}). The interval becomes
empty exactly when $H_1, \dots, H_i$ and the disc have no common point, and
then the program is infeasible (line~\ref{alg:ch13-lp:fail}). Each call of
\OrcaLPLine costs $\bigO{i}$; the total is $\bigO{m^2}$ in the worst case
and, if the half-planes are processed in random order, $\bigO{m}$ in
expectation (\cref{thm:ch13-lp}).

\subsection{Infeasibility and the dense fallback}
\label{sec:ch13-dense}

In a dense crowd the half-planes may have an empty intersection: every
velocity in the disc violates at least one of them, which means that no
velocity is safe for the full horizon $\ttc$ with respect to every
neighbour. This is not a failure of the algorithm but a fact about the
situation, and the agent still has to choose a velocity. The standard remedy,
proposed with \orca and implemented in RVO2, is to relax all agent
half-planes uniformly and choose the velocity that \textbf{minimises the
largest violation}\index{ORCA!dense fallback}\index{linear program!maximum penetration},
\begin{equation}
  \min_{\vel \in \disc{\vect{0}}{v_{\max}}}\ \max_{i}\ \mathrm{viol}_i(\vel),
  \qquad \mathrm{viol}_i(\vel) = -(\vel - \vect{q}_i)\cdot\vect{n}_i,
  \label{eq:ch13-dense}
\end{equation}
while the half-planes of static obstacles remain hard. Relaxing every agent
half-plane by the same amount $\delta$ and taking the smallest $\delta$ for
which a velocity survives is exactly \cref{eq:ch13-dense}; van den Berg et
al.\ call the result the \emph{safest possible velocity}, the one that
penetrates the half-planes least~\cite{vandenberg2011orca}. It is a
heuristic: it minimises a penetration, not the time to the first collision.
It also ignores $\vel^{\mathrm{pref}}$ entirely, because safety comes first.
\Cref{eq:ch13-dense} is a linear program in the three variables
$(v_x, v_y, \delta)$ together with the speed disc, the same
half-plane-and-disc structure as \cref{alg:ch13-lp},
and \cref{alg:ch13-dense} solves it with the two-dimensional machinery of
\cref{alg:ch13-lp}: it walks through the half-planes in order and, whenever
one of them is violated more deeply than the current largest violation,
finds the velocity that minimises the violation of that half-plane subject
to the constraint that no earlier half-plane is violated more. The
constraint ``$H_j$ is violated no more than $H_i$'' is itself a half-plane:
$\mathrm{viol}_j(\vel) \le \mathrm{viol}_i(\vel)$ is linear in $\vel$, its
boundary passes through the intersection point of the two boundary lines,
where both violations vanish, and it is the bisector of the two lines.
Minimising $\mathrm{viol}_i$ means moving as far as possible in direction
$\vect{n}_i$, which is \cref{alg:ch13-lp} with a linear objective instead of
a distance, the variant that the implementation calls with
\code{direction\_opt}.

\begin{algorithm}[htb]
\caption{Dense fallback: minimise the largest violation when the half-planes do not intersect}
\label{alg:ch13-dense}
\KwIn{half-planes $H_1, \dots, H_m$, of which $H_1, \dots, H_k$ (static obstacles) are hard; the index $f > k$ of the half-plane at which \cref{alg:ch13-lp} failed; the velocity $\vel$ it returned, which satisfies $H_1, \dots, H_{f-1}$; $v_{\max}$}
\KwOut{a velocity in the disc that satisfies the hard half-planes and minimises $\max_{i > k} \mathrm{viol}_i(\vel)$}
\Fn{\OrcaDense{$H_1, \dots, H_m$, $k$, $f$, $\vel$, $v_{\max}$}}{
  $\delta \gets 0$\tcp*{largest violation so far}
  \For{$i \gets f$ \KwTo $m$}{
    \If{$\mathrm{viol}_i(\vel) > \delta$}{\label{alg:ch13-dense:deeper}
      $\mathcal{P} \gets (H_1, \dots, H_k)$\tcp*{hard half-planes stay as they are}
      \ForEach{$j \in \set{k+1, \dots, i-1}$}{
        append to $\mathcal{P}$ the half-plane $\set{\vel : \mathrm{viol}_j(\vel) \le \mathrm{viol}_i(\vel)}$, bounded by the bisector of the boundary lines of $H_i$ and $H_j$\label{alg:ch13-dense:bisector}\;
      }
      $\vel' \gets$ the point of $\bigcap\mathcal{P} \cap \disc{\vect{0}}{v_{\max}}$ furthest in direction $\vect{n}_i$ (\OrcaLP with a linear objective)\label{alg:ch13-dense:lp}\;
      \lIf{$\vel'$ exists}{$\vel \gets \vel'$}
      $\delta \gets \mathrm{viol}_i(\vel)$\;
    }
  }
  \KwRet{$\vel$}\;
}
\end{algorithm}

After the loop, $\delta$ is the smallest achievable maximum violation and
$\vel$ attains it. If even the hard half-planes and the disc have no common
point, the agent is inside an obstacle and no velocity is safe; the
implementation then returns the best velocity for the hard constraints it
could satisfy and lets the higher layers (\cref{ch:ch24}) react. In the
simulations of this chapter the fallback is never needed
(\cref{sec:ch13-simulation}); it becomes essential in crowds of hundreds of
agents at close spacing, which is the setting the method was designed for.

\subsection{The simulation step}
\label{sec:ch13-step}

One time step for $n$ agents puts the pieces together. Every agent builds
its half-planes from the same snapshot of positions and velocities: first
the hard half-planes of the static obstacles within sensing range, with
share $1$ and a shorter horizon $\ttc_{\mathrm{obst}}$, then one half-plane
per neighbour, with share $\tfrac12$ if the neighbour runs \orca and share
$1$ if it is an intruder. It then solves \cref{alg:ch13-lp} for the velocity
closest to $\vel^{\mathrm{pref}}$ and falls back to \cref{alg:ch13-dense}
when the program is infeasible. Only when every agent has chosen does the
simulation move all of them by $\dt$ times their new velocities
(\code{simulate} in \code{ch13\_orca.py}). The order of the half-planes
matters in two places: the obstacle half-planes come first so that
\cref{alg:ch13-dense} can keep them hard, and a random order of the agent
half-planes gives the expected linear running time.

% ---------------------------------------------------------------------
\section{A worked example}
\label{sec:ch13-example}

\begin{example}[Two agents, one half-plane each]
\label{ex:ch13-two-agents}
Agent A is at $\pos_A = (0, 0)$ with velocity $\vel_A = (1, 0)$, agent B at
$\pos_B = (4, 0.5)$ with velocity $\vel_B = (-1, 0)$; both have radius
$0.5$\,m, speed limit $v_{\max} = 1.5$\,m/s and preferred velocities
$\vel_A^{\mathrm{pref}} = (1.2, 0)$ and $\vel_B^{\mathrm{pref}} = (-1.2, 0)$,
slightly faster than they currently fly. The horizon is $\ttc = 4$\,s and
the time step $\dt = 0.1$\,s. Every number below is printed by
\code{worked\_example()} in \code{ch13\_orca.py} and collected in
\cref{tab:ch13-example}; \cref{fig:ch13-orca-construction} draws the same
numbers.
\end{example}

\begin{table}[tb]
  \centering\small
  \caption{The \orca construction of \cref{ex:ch13-two-agents}, computed by
  \code{ch13\_orca.py} (four decimals). The half-plane of A is
  $-0.1260\,(v_x - 0.9841) - 0.9920\,(v_y + 0.1250) \ge 0$.}
  \label{tab:ch13-example}
  \begin{tabular}{@{}lll@{}}
  \toprule
  Step & Quantity & Value\\
  \midrule
  1 & $\pos_{\mathrm{rel}}$, $R$, $\ttc$ & $(4, 0.5)$, $1$, $4$\\
    & $\norm{\pos_{\mathrm{rel}}}$, $\varphi$, $\theta = \arcsin(R/\norm{\pos_{\mathrm{rel}}})$ & $4.0311$, $7.13^\circ$, $14.36^\circ$\\
    & truncating disc: centre $\pos_{\mathrm{rel}}/\ttc$, radius $R/\ttc$ & $(1, 0.125)$, $0.25$\\
  2 & $\vel_{\mathrm{rel}} = \vel_A - \vel_B$ & $(2, 0)$\\
  3 & $\vect{w} = \vel_{\mathrm{rel}} - \pos_{\mathrm{rel}}/\ttc$ & $(1, -0.125)$\\
    & $\vect{w}\cdot\pos_{\mathrm{rel}}$ (not $< 0$: not the arc) & $3.9375$\\
    & $\pos_{\mathrm{rel}} \times \vect{w}$ ($< 0$: right leg) & $-1$\\
    & $\ell = \sqrt{\norm{\pos_{\mathrm{rel}}}^2 - R^2}$, $\vect{d}_{\mathrm{R}}$ & $3.9051$, $(0.9920, -0.1260)$\\
    & $\vect{q} = (\vel_{\mathrm{rel}}\cdot\vect{d}_{\mathrm{R}})\,\vect{d}_{\mathrm{R}}$ & $(1.9682, -0.2500)$\\
    & $\vect{u} = \vect{q} - \vel_{\mathrm{rel}}$, $\norm{\vect{u}}$ & $(-0.0318, -0.2500)$, $0.2520$\\
    & $\vect{n} = (d_{\mathrm{R},y}, -d_{\mathrm{R},x})$ & $(-0.1260, -0.9920)$\\
  4 & A: point $\vel_A + \tfrac12\vect{u}$, normal $\vect{n}$ & $(0.9841, -0.1250)$, $(-0.1260, -0.9920)$\\
    & B: point $\vel_B - \tfrac12\vect{u}$, normal $-\vect{n}$ & $(-0.9841, 0.1250)$, $(0.1260, 0.9920)$\\
  LP & violation of $\vel_A^{\mathrm{pref}}$ & $0.1512$\\
    & $\norm{\vel_A^{\mathrm{new}} - \vel_A}$, the change of A's velocity & $0.2350$\\
    & $\vel_A^{\mathrm{new}}$, $\vel_B^{\mathrm{new}}$ & $(1.1809, -0.1500)$, $(-1.1809, 0.1500)$\\
    & new $\vel_{\mathrm{rel}}$, $\angle(\vel_{\mathrm{rel}}^{\mathrm{new}}, \pos_{\mathrm{rel}})$ & $(2.3618, -0.3000)$, $-14.36^\circ = -\theta$\\
  $\alpha = 1$ & point $\vel_A + \vect{u}$, violation of $\vel_A^{\mathrm{pref}}$ & $(0.9682, -0.2500)$, $0.2772$\\
    & $\vel_A^{\mathrm{new}}$ with full responsibility & $(1.1651, -0.2750)$\\
  \bottomrule
  \end{tabular}
\end{table}

\emph{Step 1.} $\pos_{\mathrm{rel}} = (4, 0.5)$ has length $4.0311$ and
direction $\varphi = 7.13^\circ$; the cone half-angle is $\theta =
\arcsin(1/4.0311) = 14.36^\circ$, so in the world frame the legs point at
$21.49^\circ$ (left) and $-7.24^\circ$ (right). From here on, directions are
measured from the axis of the cone: for $\vel \ne \vect{0}$ write
$\angle(\vel, \pos_{\mathrm{rel}}) = \mathrm{atan2}(v_y, v_x) - \varphi$ for the
signed angle from $\pos_{\mathrm{rel}}$ to $\vel$, positive counter-clockwise,
so that the legs are at $\pm\theta$ and the cone is the set of directions with
$\abs{\angle(\vel, \pos_{\mathrm{rel}})} < \theta$. The truncating disc has centre $(1, 0.125)$ and
radius $0.25$. \emph{Step 2.} $\vel_{\mathrm{rel}} = (2, 0)$ points along
the $x$-axis, that is
$\angle(\vel_{\mathrm{rel}}, \pos_{\mathrm{rel}}) = -7.13^\circ$, between the
legs, and at distance $2$ it is far beyond the
truncating disc: the agents collide in about $1.6$\,s, well within
$\ttc$. \emph{Step 3.} $\vect{w} = (1, -0.125)$ has
$\vect{w}\cdot\pos_{\mathrm{rel}} = 3.9375 > 0$, so the arc is ruled out;
$\pos_{\mathrm{rel}} \times \vect{w} = -1 < 0$ selects the right leg, whose
direction is $\vect{d}_{\mathrm{R}} = (0.9920, -0.1260)$. The foot of the
perpendicular is $\vect{q} = (1.9682, -0.2500)$, hence
$\vect{u} = (-0.0318, -0.2500)$ with $\norm{\vect{u}} = 0.2520$\,m/s, and the
outward normal is $\vect{n} = (-0.1260, -0.9920)$: the cheapest way out of
the cone is for the relative velocity to turn slightly downwards,
that is, for A to pass below B. \emph{Step 4.} A's half-plane has the point
$\vel_A + \tfrac12\vect{u} = (0.9841, -0.1250)$ and the normal $\vect{n}$;
B's half-plane has the point $\vel_B - \tfrac12\vect{u} = (-0.9841, 0.1250)$
and the normal $-\vect{n}$.

\emph{The program.} A's preferred velocity $(1.2, 0)$ violates its
half-plane by $0.1512$: it is inside the disc, so \cref{alg:ch13-lp} projects
it onto the boundary line and returns
$\vel_A^{\mathrm{new}} = (1.2, 0) + 0.1512\,\vect{n} = (1.1809, -0.1500)$;
by symmetry B returns $(-1.1809, 0.1500)$. Each agent ends up $0.1512$\,m/s
away from its preferred velocity (and $0.2350$\,m/s away from the velocity it
currently flies; the half-plane only required
$\tfrac12\norm{\vect{u}} = 0.1260$\,m/s along $\vect{n}$), and the new relative
velocity $(2.3618, -0.3000)$ has
$\angle(\vel_{\mathrm{rel}}^{\mathrm{new}}, \pos_{\mathrm{rel}}) = -14.36^\circ = -\theta$:
it lies exactly on the right leg, so the discs will graze and not
penetrate. This is the pairwise guarantee of
\cref{thm:ch13-pairwise} in numbers. \emph{Full responsibility.} If B were an
intruder that keeps flying $(-1, 0)$, A would use $\alpha = 1$: the point
moves to $(0.9682, -0.2500)$, the violation of the preferred velocity
doubles to $0.2772$, and A alone chooses $(1.1651, -0.2750)$, whose
relative velocity against the unchanged $\vel_B$ again lies on the leg. Had
A used $\alpha = \tfrac12$ against a non-reacting B, the relative velocity
would have been $(2.1809, -0.1500)$, at $-11.06^\circ$ from
$\pos_{\mathrm{rel}}$ ($-3.93^\circ$ in the world frame), inside the cone of
half-angle $14.36^\circ$.

% ---------------------------------------------------------------------
\section{Eight agents on a circle}
\label{sec:ch13-simulation}

The classic test of a reciprocal method puts $n$ agents on a circle and
sends each to the antipodal point, so that all straight paths cross at the
centre at the same time. \Cref{fig:ch13-circle} shows eight agents of radius
$0.5$\,m on a circle of radius $10$\,m, cruise speed $1$\,m/s, speed limit
$1.5$\,m/s, $\ttc = 5$\,s and $\dt = 0.1$\,s, simulated for $30$\,s with
\code{simulate()} of \code{ch13\_orca.py} in four modes: no avoidance;
the sampled velocity obstacle rule (the velocity closest to
$\vel^{\mathrm{pref}}$ among $721$ candidates on a polar grid whose time to
collision with every neighbour is at least $\ttc$); sampled RVO (the same
with the reciprocal test $2\vel - \vel_A$); and \orca. The preferred
velocities are perturbed by fixed offsets of $0.02$\,m/s so that the
instance is not perfectly symmetric, as no real instance is.
\Cref{tab:ch13-circle} gives the numbers, all printed by
\code{gen\_ch13\_circle.py}.

\begin{figure}[tb]
  \centering
  \inputfigure{ch13/circle}
  \caption{Eight agents swapping antipodal positions on a circle. Top: the
  traces (circles mark the starts) with no avoidance, with the sampled
  velocity obstacle rule and with \orca. Bottom: the smallest centre
  distance over all pairs against time, for all four methods (the lower panel
  also shows sampled RVO, which gets no trace panel); the dashed line is the collision
  threshold $r_A + r_B = 1$\,m. Without avoidance all agents meet at the
  centre; the VO agents zig-zag and four pairs still collide; the \orca
  agents spiral gently around the centre and the smallest distance touches
  $1.000$\,m without ever going below it.}
  \label{fig:ch13-circle}
\end{figure}

\begin{table}[tb]
  \centering\small
  \caption{Results of the circle scenario of \cref{fig:ch13-circle}
  ($30$\,s, $300$ steps). ``Colliding pairs'' counts the pairs whose discs
  overlapped at some step; ``reversals'' counts the steps at which the
  velocity change of agent~0 turned against the previous change, and
  ``velocity variation'' is the sum of its velocity changes over the run.}
  \label{tab:ch13-circle}
  \setlength{\tabcolsep}{4.5pt}
  \begin{tabular}{@{}lrrrrrrr@{}}
  \toprule
  Method & \makecell[r]{colliding\\steps} & \makecell[r]{colliding\\pairs} & \makecell[r]{min.\ sep.\\(m), at $t$ (s)} & \makecell[r]{last\\arrival (s)} & \makecell[r]{mean path\\length (m)} & \makecell[r]{reversals\\agent 0} & \makecell[r]{velocity\\var.\ (m/s)}\\
  \midrule
  no avoidance & 33 & 28 & $0.010$, $9.9$ & $20.3$ & $20.00$ & 1 & $2.05$\\
  sampled VO   & 25 &  4 & $0.393$, $11.0$ & $23.6$ & $20.45$ & 62 & $37.82$\\
  sampled RVO  &  0 &  0 & $1.002$, $14.2$ & $23.5$ & $20.54$ & 33 & $15.13$\\
  \orca        &  0 &  0 & $1.000$, $14.8$ & $26.2$ & $20.29$ & 2 & $3.37$\\
  \bottomrule
  \end{tabular}
\end{table}

Three things to notice. Without avoidance every pair collides at the centre
at $t \approx 10$\,s, as it must. The velocity obstacle rule avoids most of
that, but the dance of \cref{thm:ch13-dance} now involves seven neighbours
at once: agent~0 reverses its velocity change $62$ times and its velocity
changes add up to $37.8$\,m/s over the run, and the stale-velocity problem is
severe enough that four pairs still overlap, with a smallest separation of
$0.39$\,m against a threshold of $1$\,m. RVO stays collision-free here, but
it still reverses $33$ times: the reciprocal dance among many agents, which
happens to be harmless in this instance and is not guaranteed to be. \orca
is collision-free with two reversals, a velocity variation of $3.4$\,m/s and
paths only $1.5\,\%$ longer than the straight line; the linear program was
feasible at every one of the $2\,400$ agent-steps. Its price is time: the
agents slow down near the centre and the last one arrives at $26.2$\,s
instead of $20$\,s, because a half-plane never asks an agent to speed up
past $v_{\max}$ and the projections of the preferred velocities point
partly backwards while the cluster passes. That slowing-down is the
behaviour the local safety layer of \cref{ch:ch24} wants: it buys time for
the global planner instead of inventing a detour.

% ---------------------------------------------------------------------
\section{Properties: what \orca guarantees and what it does not}
\label{sec:ch13-properties}

\subsection{The pairwise guarantee}

\begin{lemma}[The truncated velocity obstacle is convex]
\label{thm:ch13-convex}
For $\norm{\pos_{\mathrm{rel}}} > R$ the set $\VO^{\ttc}_{A|B}$ is open and
convex, and at every boundary point $\vect{q}$ with outward unit normal
$\vect{n}$ it lies entirely on the inner side of the tangent line:
$\VO^{\ttc}_{A|B} \subseteq \set{\vel : (\vel - \vect{q})\cdot\vect{n} < 0}$.
\end{lemma}
\begin{proof}
Openness follows from \cref{eq:ch13-vo-tau}, a union of open discs. For
convexity let $\vel_1, \vel_2 \in \VO^{\ttc}_{A|B}$ with times
$t_1, t_2 \in (0, \ttc]$ such that $t_i\vel_i \in \disc{\pos_{\mathrm{rel}}}{R}$,
let $\lambda \in [0, 1]$ and $\vel = \lambda\vel_1 + (1-\lambda)\vel_2$.
Define $t$ by $1/t = \lambda/t_1 + (1-\lambda)/t_2$; then
$t \le \max(t_1, t_2) \le \ttc$ and
\[
  t\,\vel = \frac{t\lambda}{t_1}\,(t_1\vel_1) + \frac{t(1-\lambda)}{t_2}\,(t_2\vel_2)
\]
is a convex combination of two points of the disc, because the two
coefficients are non-negative and sum to one. Hence $t\vel$ lies in the
disc and $\vel \in \VO^{\ttc}_{A|B}$. The boundary of a convex set with a
tangent line at $\vect{q}$ is supported by that line, and the boundary of
$\VO^{\ttc}_{A|B}$ has a tangent line everywhere: on the arc and on the legs
it is smooth, and at the two points where they meet the leg is tangent to
the circle. Openness turns the weak inequality of the supporting line into a
strict one.
\end{proof}

\begin{theorem}[Pairwise collision avoidance]
\label{thm:ch13-pairwise}
Let agents $A$ and $B$ with $\norm{\pos_{\mathrm{rel}}} > R$ know each
other's positions, radii and current velocities exactly, and let both
build their half-planes with the same horizon $\ttc$ from the same current
velocities. If $\vel_A^{\mathrm{new}} \in \ORCA^{\ttc}_{A|B}$ and
$\vel_B^{\mathrm{new}} \in \ORCA^{\ttc}_{B|A}$, then
$\vel_A^{\mathrm{new}} - \vel_B^{\mathrm{new}} \notin \VO^{\ttc}_{A|B}$:
flying these velocities, the discs of $A$ and $B$ do not overlap during the
next $\ttc$ seconds.
\end{theorem}
\begin{proof}
By \cref{eq:ch13-orca,eq:ch13-orca-mirror},
$(\vel_A^{\mathrm{new}} - \vel_A - \tfrac12\vect{u})\cdot\vect{n} \ge 0$ and
$(\vel_B^{\mathrm{new}} - \vel_B + \tfrac12\vect{u})\cdot(-\vect{n}) \ge 0$.
Adding the two inequalities gives
\[
  \bigl(\vel_A^{\mathrm{new}} - \vel_B^{\mathrm{new}} - (\vel_A - \vel_B) - \vect{u}\bigr)\cdot\vect{n}
  = \bigl(\vel_{\mathrm{rel}}^{\mathrm{new}} - \vect{q}\bigr)\cdot\vect{n} \ge 0,
\]
where $\vect{q} = \vel_{\mathrm{rel}} + \vect{u}$ is the boundary point of
Step~3 and $\vect{n}$ its outward normal. By \cref{thm:ch13-convex} every
point of $\VO^{\ttc}_{A|B}$ satisfies the opposite strict inequality, so
$\vel_{\mathrm{rel}}^{\mathrm{new}} \notin \VO^{\ttc}_{A|B}$, and by
\cref{def:ch13-vo} the discs do not overlap for $t \in (0, \ttc]$.
\end{proof}

The proof uses nothing about how the two agents chose their velocities
inside their half-planes, which is what makes the guarantee compose. The
same computation with shares $\alpha_A$, $\alpha_B$ produces
$(\alpha_A + \alpha_B)\,\vect{u}$ in place of $\vect{u}$, and the conclusion
holds exactly when $\alpha_A + \alpha_B = 1$ (\cref{exr:ch13-shares}); with
$\alpha_B = 0$ the agent $B$ contributes nothing and must therefore keep the
velocity $\vel_B$ that $A$ used, which is the non-cooperative case of
\cref{sec:ch13-noncooperative}.

\begin{corollary}[Collision-free motion of $n$ agents]
\label{thm:ch13-nbody}
Let $n$ agents run the step of \cref{sec:ch13-step} with the same horizon $\ttc$ from
the same snapshot, with share $\tfrac12$ towards each other and share $1$
towards obstacles and intruders that keep their assumed velocities. If every
agent's program \cref{eq:ch13-program} is feasible, then no two agents and
no agent and obstacle overlap during the next $\ttc$ seconds. If the step is
repeated every $\dt \le \ttc$ and the programs stay feasible, the motion is
collision-free for as long as this holds.
\end{corollary}
\begin{proof}
A feasible program returns a velocity in every half-plane of the agent
(\cref{thm:ch13-lp}), so \cref{thm:ch13-pairwise} applies to every pair
of agents and its non-cooperative variant to every agent--obstacle pair.
Each guarantee covers $\ttc \ge \dt$ seconds, hence the whole step, and
the next step starts from the new snapshot.\index{ORCA!guarantee}
\end{proof}

\subsection{The program}

\begin{proposition}[The per-agent program]
\label{thm:ch13-lp}
If the feasible set of \cref{eq:ch13-program} is non-empty, the program has
a unique solution and \cref{alg:ch13-lp} returns it after $\bigO{m^2}$
operations in the worst case; if the half-planes are processed in uniformly
random order, the expected number of operations is $\bigO{m}$. If the
feasible set is empty, the algorithm reports the smallest index $i$ such
that $H_1 \cap \dots \cap H_i \cap \disc{\vect{0}}{v_{\max}} = \emptyset$.
\end{proposition}
\begin{proof}[Proof sketch]
The feasible set is closed and convex and the objective is strictly convex,
which gives existence and uniqueness. Correctness is the invariant of the
walkthrough: the optimum for $H_1, \dots, H_i$ equals the optimum for
$H_1, \dots, H_{i-1}$ if the latter satisfies $H_i$, and lies on the
boundary line of $H_i$ otherwise, where \OrcaLPLine finds it exactly by a
one-dimensional minimisation over an interval. For the running time,
iteration $i$ costs $\bigO{1}$ if $H_i$ is satisfied and $\bigO{i}$
otherwise. Backwards analysis~\cite{seidel1991small,deberg2008computational}:
the optimum for $H_1, \dots, H_i$ is determined by at most two of the
half-planes (its boundary point lies on at most two boundary lines, or on one
line and the circle), so if the order is random the probability that $H_i$
is one of the half-planes that determine it, which is the only case in
which $H_i$ can be violated by the optimum for the first $i-1$, is at most
$2/i$. The expected cost of iteration $i$ is $\bigO{1} + \bigO{i}\cdot 2/i =
\bigO{1}$, and the total is $\bigO{m}$.
\end{proof}

With $k$ neighbours per agent, one step of \cref{sec:ch13-step} costs
$\bigO{k}$ per agent for the half-planes and the program, plus the
neighbour query, $\bigO{\log n + k}$ with a kd-tree, so $\bigO{n(\log n + k)}$
per step for the whole swarm. The RVO2 library simulates thousands of agents
in real time on one processor core.\index{ORCA!complexity}

\subsection{What \orca does not guarantee}
\label{sec:ch13-limits}

The guarantee of \cref{thm:ch13-nbody} is conditional, and every condition
fails somewhere in a real swarm; the honest summary is that \orca is a
\emph{safety filter} for the next $\ttc$ seconds and nothing more. The
half-planes may not intersect, and the fallback of \cref{alg:ch13-dense}
then returns the least-penetrating velocity, which does not exclude a
collision. \Cref{thm:ch13-pairwise} also assumes $\norm{\pos_{\mathrm{rel}}} > R$:
once two discs already overlap, the overlap branch of \cref{thm:ch13-cases}
carries no $\ttc$-guarantee at all and only promises that the pair separates
after one time step $\dt$. Both agents must
use the same $\pos_{\mathrm{rel}}$, $R$, $\ttc$ and current velocities;
sensing noise, latency and different horizons make the two half-planes
slightly inconsistent, which the radius margin of
\cref{sec:ch13-implementation} absorbs, while the intruder case absorbs the
large inconsistencies by taking full responsibility. The agents are assumed
to fly the new velocity immediately, whereas a drone tracks it with a delay
and an error (\cref{ch:ch02,ch:ch14}); again the margin pays for it.
Finally, \orca chooses the closest permitted velocity to the preferred one,
step by step, with no memory and no lookahead beyond $\ttc$, so it can stop
an agent forever.

\begin{proposition}[\orca is not a global planner]
\label{thm:ch13-incomplete}
There are instances with a collision-free solution in which agents running the step of \cref{sec:ch13-step} never reach their goals.
\end{proposition}
\begin{proof}[Proof by example]
Two agents of radius $r_0$ face each other in a corridor of width $3r_0$
with a side bay near one end, so that they can pass only if one agent
enters the bay and waits. Both preferred velocities point along the
corridor. The wall half-planes bound the lateral velocity by a fixed multiple
of the forward velocity (\cref{exr:ch13-corridor} computes the factor $0.75$
for one geometry), so a lateral escape is available only while the agent still
moves forward. As the two agents approach head-on, the arc case of
\cref{thm:ch13-cases} applies: with $\hat{\pos} = \pos_{\mathrm{rel}}/\norm{\pos_{\mathrm{rel}}}$
and the gap $d = \norm{\pos_{\mathrm{rel}}} - R$ it gives
$\vect{u} = (d/\ttc)\,\hat{\pos} - \vel_{\mathrm{rel}}$ and
$\vect{n} = -\hat{\pos}$, and the half-plane through
$\vel_A + \tfrac12\vect{u}$ therefore permits each agent to approach the other
at no more than $d/(2\ttc)$, whatever its current speed. The gap obeys
$d_{k+1} = d_k(1 - \dt/\ttc)$ and both speeds decay geometrically to zero.
Backing into the bay is further from $\vel^{\mathrm{pref}}$ than slowing down
and no half-plane ever asks for it, so neither agent ever reaches its goal,
although a collision-free joint plan exists.
\end{proof}

Perfectly symmetric configurations produce a milder version of the same
effect, agents that slow down and hover instead of passing, which is why
implementations break symmetry by a small random perturbation of the
preferred velocities, as \cref{sec:ch13-simulation} did. The remedy for
deadlocks is not inside \orca: it is the global planner of the hybrid
architecture, which sequences the corridor in the nominal paths, and the
replanning layer, which detects that an agent has been stopped for too long
and plans around (\cref{ch:ch24}).\index{ORCA!deadlock}\index{completeness!ORCA}

% ---------------------------------------------------------------------
\section{Variants and extensions}
\label{sec:ch13-variants}

\paragraph{HRVO.}
The hybrid reciprocal velocity obstacle~\cite{snape2011hrvo} attacks the
side-choice problem of RVO directly: for the side on which the agent is
already passing it uses the RVO leg, for the other side the VO leg, so that
crossing over to the wrong side requires the full change instead of half.
It reduces the reciprocal dance markedly on real differential-drive robots,
but it keeps the non-convex cone and therefore the sampled velocity choice
and the lack of an $n$-agent guarantee.

\paragraph{Static obstacles.}
A static disc is an agent with $\vel_B = \vect{0}$ and share $1$. For
polygonal obstacles the velocity obstacle of a line segment is the union of
the velocity obstacles of its points, whose boundary consists of two legs
tangent to the discs around the endpoints and, after truncation, of a
stadium-shaped front; RVO2 builds one half-plane per visible segment with
full responsibility and a separate, shorter horizon $\ttc_{\mathrm{obst}}$,
and places these half-planes first so that they stay hard in the dense
fallback.

\paragraph{Non-holonomic robots.}
A differential-drive robot cannot fly an arbitrary velocity instantly.
Alonso-Mora et al.\ keep the \orca half-planes and restrict the velocity
choice to the set of holonomic velocities that the robot can track within a
bounded error $\eta$ over the next interval, a non-convex set that they
approximate by convex polygons, while inflating the radius by
$\eta$~\cite{alonsomora2013optimal}. The same recipe applies to a
multirotor with a velocity controller: measure the tracking error, inflate
the radius by it, and clip the velocity change per step to what the
controller delivers (\cref{ch:ch14} treats the acceleration limits
explicitly).

\paragraph{Uncertainty.}
An intruder's future velocity is a prediction with a covariance
(\cref{ch:ch18,ch:ch20}). Inflating $r_B$ by a multiple of the predicted
position standard deviation at the horizon turns \cref{eq:ch13-orca-full}
into a conservative constraint, which is what the drone system of this book
does. Zhu and Alonso-Mora formulate the same idea as a chance constraint,
a collision probability below a chosen level, inside a model predictive
controller for multirotors~\cite{zhu2019chance}; \cref{ch:ch21} introduces
that controller, and the linearised chance constraint has exactly the shape
of a half-plane with an uncertainty-dependent margin.

% ---------------------------------------------------------------------
\section{Implementation notes}
\label{sec:ch13-implementation}

The file \code{ch13\_orca.py} implements everything in this chapter with
plain tuples and a few NumPy calls; \cref{lst:ch13-cases,lst:ch13-lp}
show the functions that carry the mathematics, where the code writes the
combined radius \code{r} for what the text calls $R$, and the line program of
\cref{alg:ch13-lp} is \code{\_lp\_on\_line} in the same file. The geometry helpers
(\code{dot}, \code{cross}, \code{ray\_circle\_intersection}, \dots) are
copied from \code{ch02\_toolbox.py} so that the file stands alone.

\begin{lstlisting}[caption={The case analysis of \cref{thm:ch13-cases} (from \code{code/ch13\_orca.py}, docstring omitted). The function returns $\vect{u}$, $\vect{n}$ and the name of the case.},label={lst:ch13-cases}]
def vo_closest_boundary_point(p, v_rel, r, tau, dt):
    dist2 = dot(p, p)
    r2 = r * r
    if dist2 > r2:
        c = scale(p, 1.0 / tau)              # centre of the truncating disc
        rho = r / tau                        # its radius
        w = sub(v_rel, c)                    # from that centre to v_rel
        w2 = dot(w, w)
        wp = dot(w, p)
        if wp < 0.0 and wp * wp > r2 * w2:
            # angle(w, -p) < arccos(r / |p|): the closest point is on the arc
            w_len = math.sqrt(w2)
            n = scale(w, 1.0 / w_len)
            return scale(n, rho - w_len), n, "disc"
        leg = math.sqrt(dist2 - r2)          # length of a tangent from 0
        if cross(p, w) > 0.0:
            d = ((p[0] * leg - p[1] * r) / dist2,
                 (p[0] * r + p[1] * leg) / dist2)   # unit vector, left leg
            n = (-d[1], d[0])                        # outward: d rotated +90
            case = "left leg"
        else:
            d = ((p[0] * leg + p[1] * r) / dist2,
                 (-p[0] * r + p[1] * leg) / dist2)  # unit vector, right leg
            n = (d[1], -d[0])                        # outward: d rotated -90
            case = "right leg"
        q = scale(d, dot(v_rel, d))              # projection onto the leg
        return sub(q, v_rel), n, case
    # the discs already overlap: leave the disc of centre p/dt, radius r/dt
    c = scale(p, 1.0 / dt)
    rho = r / dt
    w = sub(v_rel, c)
    w_len = norm(w)
    if w_len > EPS:
        n = scale(w, 1.0 / w_len)
    elif dist2 > 0.0:
        n = unit(scale(p, -1.0))
    else:
        n = (1.0, 0.0)
    return scale(n, rho - w_len), n, "overlap"
\end{lstlisting}

\begin{lstlisting}[caption={The half-plane (\cref{alg:ch13-halfplane}) and the incremental program (\cref{alg:ch13-lp}); \code{HalfPlane.violation(v)} is $-(\vel - \vect{q})\cdot\vect{n}$ and \code{direction} is $\vect{n}$ rotated by $-90^\circ$.},label={lst:ch13-lp}]
def orca_half_plane(agent, other, tau, dt=0.1, reciprocal=True):
    p = sub(other.position, agent.position)
    v_rel = sub(agent.velocity, other.velocity)
    r = agent.radius + other.radius
    u, n, case = vo_closest_boundary_point(p, v_rel, r, tau, dt)
    share = 0.5 if reciprocal else 1.0
    return OrcaLine(n, add(agent.velocity, scale(u, share)), u, case)


def lp_incremental(lines, v_max, v_opt, direction_opt=False):
    if direction_opt:
        v = scale(v_opt, v_max)
    elif dot(v_opt, v_opt) > v_max * v_max:
        v = scale(v_opt, v_max / norm(v_opt))
    else:
        v = v_opt
    for i, line in enumerate(lines):
        if line.violation(v) > 0.0:
            new_v = _lp_on_line(lines, i, v_max, v_opt, direction_opt)
            if new_v is None:
                return v, i
            v = new_v
    return v, len(lines)
\end{lstlisting}

\paragraph{Libraries.}
The reference implementation is the RVO2 library of the \orca authors
(C++, with ports to several languages, and RVO2-3D for half-spaces). Its
per-agent parameters map directly onto this chapter: \code{radius},
\code{maxSpeed}, \code{timeHorizon} ($\ttc$), \code{timeHorizonObst}
($\ttc_{\mathrm{obst}}$), \code{neighborDist} and \code{maxNeighbors} (the
sensing range and the cap on the number of half-planes), and the global
\code{timeStep} ($\dt$). Reading its \code{Agent.cpp} after this chapter is
a good exercise: \cref{alg:ch13-halfplane,alg:ch13-lp,alg:ch13-dense} are
its \code{computeNewVelocity}, \code{linearProgram2} with
\code{linearProgram1}, and \code{linearProgram3}.

\paragraph{Neighbour queries.}
Every agent needs its neighbours within the sensing range, and a linear
scan costs $\bigO{n}$ per agent. RVO2 rebuilds a kd-tree over the agent
positions at every step and queries it; the same structure serves the
nearest-neighbour searches of \cref{ch:ch16}. Cap the number of neighbours
by distance: the closest $10$ to $20$ agents determine the program, the
rest add half-planes that are never active.

\paragraph{Horizon and time step.}
The guarantee covers $\ttc$ seconds and the velocity is recomputed every
$\dt$, so $\ttc$ must exceed $\dt$ by a comfortable factor, typically
$\ttc = 10$ to $50\,\dt$. A larger $\ttc$ makes agents react earlier and more
gently, but it also makes the half-planes more restrictive, since a
half-plane forbids velocities that lead to a collision anywhere within
$\ttc$, and it makes infeasibility and hovering more likely in crowds.
Obstacles get a shorter $\ttc_{\mathrm{obst}}$\index{time horizon!obstacles}, because a wall is not
to move out of the way and an agent that fears it from far away hugs the
middle of every corridor.

\paragraph{The preferred velocity.}
$\vel^{\mathrm{pref}}$ is the interface to the global planner: the vector
towards the next waypoint of the nominal path at cruise speed, scaled down
near the goal so that the agent does not overshoot in the last step
(\code{preferred\_velocity} in the code). It is never the raw direction to
the final goal in a cluttered map, which would make \orca a potential-field
method with all its local minima (\cref{ch:ch15}).

\paragraph{Radius inflation.}
The radius that goes into \cref{alg:ch13-halfplane} is not the physical
radius. Add the velocity tracking error of the controller over one step,
the localisation error, half of the position uncertainty of the neighbour
and, for an intruder, the predicted position uncertainty at the horizon.
The margin also converts the grazing contact that the guarantee permits into
a real clearance, and it absorbs the asynchrony between agents that do not
update at exactly the same instant.

\begin{pitfall}[Mixing reciprocal and full responsibility]
The shares must sum to one \emph{per pair}, and both sides must agree. Two
common bugs: treating an intruder that does not react as a reciprocal
neighbour ($\alpha = \tfrac12$ on one side, $0$ on the other), after which the
relative velocity moves only halfway out of the cone and the collision
course remains, as in the last lines of \cref{tab:ch13-example}; and
treating cooperative neighbours with full responsibility on both sides
($\alpha = 1$ and $1$), which is safe but is exactly the velocity obstacle
rule of \cref{thm:ch13-dance}, and the dance returns. Decide the share from
the \emph{type} of the neighbour (swarm mate, intruder, obstacle), never
from its current behaviour.
\end{pitfall}

\begin{pitfall}[The normal points the wrong way, or the disc is forgotten]
$\vect{n}$ must be the \emph{outward} normal of $\VO^{\ttc}$ at the closest
point in every case: away from the truncating disc's centre on the arc,
and the leg direction rotated by $+90^\circ$ on the left leg and by
$-90^\circ$ on the right leg. A sign error in one case gives a half-plane
on the wrong side, and the program cheerfully returns a velocity that leads
into the cone. Test every case against
\cref{thm:ch13-cases} the way the self-test does: the point
$\vel_{\mathrm{rel}} + \vect{u}$ must lie on the boundary, and
$\vel_{\mathrm{rel}} + \vect{u} \pm \eps\vect{n}$ must lie outside and
inside. The second bug is to solve the program without the disc: the line
program then has an unbounded interval, the returned velocity can exceed
$v_{\max}$, and the agent cannot fly it, so the relative velocity that the
guarantee is about is never realised.
\end{pitfall}

\begin{pitfall}[Horizon, time step and agents that must hold a formation]
With $\ttc < \dt$ the guarantee expires before the next update and agents
can overlap within one step; with $\ttc$ far larger than the time an agent
needs to cross the sensing range, every visible neighbour produces an
active half-plane and the swarm hovers or the program turns infeasible.
Choose $\ttc$ from the closing speed and the sensing range, and keep
$\ttc_{\mathrm{obst}}$ shorter. Finally, do not run \orca between the
members of a tight formation with the formation spacing as their separation:
the half-planes are permanently active, they fight the formation controller
of \cref{ch:ch23}, and the formation dissolves into hovering. Keep formation
mates apart by the controller, use \orca between formations and against
intruders, or give the leader the right of way with shares $(1, 0)$.
\end{pitfall}

% ---------------------------------------------------------------------
\section{Where this fits in the drone system}
\label{sec:ch13-drone}

\begin{dronebox}
\orca is the local safety layer of the four-layer architecture of
\cref{ch:ch24}, the Week~6 topic of the study plan (\cref{ch:appA}). It
receives from the global planner (conflict-based search \cbs, or its
bounded-suboptimal variant \ecbs, \cref{ch:ch09,ch:ch10}) the
nominal time-indexed path of each drone and turns it into the preferred
velocity: towards the next waypoint at cruise speed. While no predicted
conflict is inside the safety horizon, the half-planes are inactive and the
drone flies its nominal path exactly. When the prediction layer
(\cref{ch:ch18,ch:ch20}) reports an intruder whose predicted trajectory
enters the safety horizon, the trigger of \cref{ch:ch24} switches the local
layer on: the swarm mates within sensing range become reciprocal
half-planes with share $\tfrac12$, and the intruder becomes a half-plane
with \emph{full} responsibility, built from its predicted velocity and a
radius inflated by the predicted position uncertainty at the horizon. The layer delivers one velocity
command per step to the flight controller and reports two things upwards:
whether its program was feasible, and how far the drone has drifted from
the nominal path. When the conflict clears, the preferred velocity points at
a waypoint further along the nominal path and \orca reconnects to it by
itself. When the drift is too large, the program has been infeasible or the
drone has been stopped for longer than a threshold, the replanning layer
(\dstarlite or \astar, \cref{ch:ch05}) takes over, because \orca cannot get
out of a deadlock on its own (\cref{thm:ch13-incomplete}). The guarantee that
survives the integration is the pairwise one of \cref{thm:ch13-pairwise},
per step and per horizon; everything beyond it, reaching the goal, keeping
the formation, staying within communication range, is the job of the other
layers.
\end{dronebox}

% ---------------------------------------------------------------------
\section{Summary}
\begin{summary}
\begin{itemize}
  \item Two agents that avoid each other with plain velocity obstacles
        react to each other's stale velocities and oscillate with period two:
        both swerve, both return, both swerve.
  \item RVO makes each agent take half of the velocity change,
        $\RVO_{A|B}(\vel_B, \vel_A) = \set{\vel : 2\vel - \vel_A \in \vel_B + \VO_{A|B}}$,
        the cone with its apex moved to $\tfrac12(\vel_A + \vel_B)$. For two
        agents that pass on the same side this removes the oscillation; with
        several agents or disagreeing sides the reciprocal dance returns, and
        there is no guarantee.
  \item \orca builds, for each pair and horizon $\ttc$, the truncated cone
        $\VO^{\ttc}_{A|B}$, takes the current relative velocity, computes the
        smallest change $\vect{u}$ to the boundary (arc, left leg, right leg,
        or the one-step disc when already overlapping) with the outward normal
        $\vect{n}$, and permits $A$ the half-plane
        $\set{\vel : (\vel - (\vel_A + \tfrac12\vect{u}))\cdot\vect{n} \ge 0}$;
        $B$ gets the mirror half-plane with $-\vect{u}$ and $-\vect{n}$.
  \item Each agent chooses the velocity closest to its preferred one inside
        all its half-planes and the disc $\norm{\vel} \le v_{\max}$: a
        two-dimensional convex program solved exactly by the incremental
        linear program in expected linear time. If the intersection is empty,
        the dense fallback minimises the largest violation.
  \item Against an obstacle or an intruder that does not react, the agent
        takes full responsibility: the half-plane through $\vel_A + \vect{u}$.
        Shares must sum to one per pair. In three dimensions the same
        construction in the plane of $\pos_{\mathrm{rel}}$ and
        $\vel_{\mathrm{rel}}$ gives a half-space.
  \item If all agents use the same horizon from the same snapshot and every
        program is feasible, no pair overlaps within $\ttc$. \orca is not a
        global planner: it has no lookahead, no memory and no completeness,
        and it deadlocks in corridors and symmetric configurations.
  \item In the circle scenario, no avoidance collides for all $28$ pairs,
        the velocity obstacle rule still collides for four and reverses its
        velocity change $62$ times, RVO is collision-free but jittery, and
        \orca is collision-free with two reversals and paths $1.5\,\%$ longer
        than straight lines.
\end{itemize}
\end{summary}

\paragraph{Further reading.}
Velocity obstacles are due to Fiorini and Shiller~\cite{fiorini1998motion};
the reciprocal velocity obstacle to van den Berg, Lin and
Manocha~\cite{vandenberg2008rvo}, and \orca, with the maximality argument, the
dense fallback and the RVO2 library, to van den Berg, Guy, Lin and
Manocha~\cite{vandenberg2011orca}. The hybrid reciprocal velocity obstacle
is by Snape, van den Berg, Guy and Manocha~\cite{snape2011hrvo}, the
non-holonomic extension by Alonso-Mora et al.~\cite{alonsomora2013optimal},
and the chance-constrained formulation for multirotors by Zhu and
Alonso-Mora~\cite{zhu2019chance}. The randomised incremental linear program
is Seidel's algorithm~\cite{seidel1991small}; the textbook treatment,
including the three-dimensional version and the backwards analysis, is
Chapter~4 of de Berg, Cheong, van Kreveld and
Overmars~\cite{deberg2008computational}. LaValle's
book~\cite{lavalle2006planning} places velocity-space planning among
moving obstacles in the wider context of planning in time-varying
environments.

% ---------------------------------------------------------------------
\section{Exercises}
\label{sec:ch13-exercises}

\begin{exercise}[\difficulty{1}]
\label{exr:ch13-mirror}
Starting from \cref{def:ch13-orca}, show in detail that the half-plane of $B$
with respect to $A$ has the point $\vel_B - \tfrac12\vect{u}$ and the normal
$-\vect{n}$: write down $\VO^{\ttc}_{B|A}$, its closest boundary point to
$\vel_B - \vel_A$, and the corresponding change and normal. Then add the two
half-plane inequalities and show that the sum says that the new relative
velocity lies on the outer side of the tangent line at $\vect{q}$.
\end{exercise}

\begin{exercise}[\difficulty{1}]
\label{exr:ch13-by-hand}
Repeat the worked example of \cref{sec:ch13-example} with B at
$\pos_B = (4, -0.5)$ instead of $(4, 0.5)$, everything else unchanged. Which
leg is closest? Give $\vect{u}$, $\vect{n}$, the point and normal of A's
half-plane and $\vel_A^{\mathrm{new}}$, and explain why the numbers are the
mirror images of \cref{tab:ch13-example}. Check with
\code{orca\_half\_plane}.
\end{exercise}

\begin{exercise}[\difficulty{2}]
\label{exr:ch13-arc-case}
With $\pos_{\mathrm{rel}} = (4, 0.5)$, $R = 1$ and $\ttc = 4$ as in the
worked example, take $\vel_{\mathrm{rel}} = (0.4, 0.05)$. (a) Show that this
relative velocity is outside $\VO^{\ttc}_{A|B}$ and compute the time at which
the discs would touch. (b) Apply \cref{thm:ch13-cases}: which case applies,
and what are $\vect{q}$, $\vect{u}$ and $\vect{n}$? What does the sign of
$\vect{u}\cdot\vect{n}$ tell you? (c) Prove the claim in the proof sketch of
\cref{thm:ch13-cases} that in the leg case the foot of the perpendicular
lies beyond the tangent point, that is, on the boundary of
$\VO^{\ttc}_{A|B}$ and not on the extension of the leg towards the apex.
\end{exercise}

\begin{exercise}[\difficulty{2}]
\label{exr:ch13-shares}
Let $A$ use the share $\alpha_A$ and $B$ the share $\alpha_B$ in
\cref{alg:ch13-halfplane}. (a) Redo the proof of \cref{thm:ch13-pairwise} and
show that the new relative velocity is guaranteed to lie outside
$\VO^{\ttc}_{A|B}$ exactly when $\alpha_A + \alpha_B = 1$; explain why both
$\alpha_A + \alpha_B > 1$ and $< 1$ can fail, using the sign of
$\vect{u}\cdot\vect{n}$. (b) In the worked example, let B be an intruder that
keeps $\vel_B = (-1, 0)$ while A uses $\alpha_A = \tfrac12$. Compute the new
relative velocity and its angle from $\pos_{\mathrm{rel}}$ and conclude that
the pair is still on a collision course. (c) What happens if two cooperative
agents both use $\alpha = 1$? Relate the answer to \cref{thm:ch13-dance}.
\end{exercise}

\begin{exercise}[\difficulty{2}]
\label{exr:ch13-rvo-truncated}
Prove \cref{thm:ch13-rvo-apex} for the truncated cone: show that
$\set{\vel : 2\vel - \vel_A \in \vel_B + \VO^{\ttc}_{A|B}} =
\tfrac12(\vel_A + \vel_B) + \VO^{2\ttc}_{A|B}$. Why does the horizon double,
and what does this mean for an implementation that wants RVO agents to react
to collisions within $\ttc$ seconds?
\end{exercise}

\begin{exercise}[\difficulty{2}]
\label{exr:ch13-lp-by-hand}
Take the three half-planes of \cref{fig:ch13-lp}, $H_1$ through
$(0.9, -0.1)$ with normal $(-0.2, -0.98)$, $H_2$ through $(-0.4, 0.2)$ with
normal $(0.94, 0.34)$, $H_3$ through $(0.5, -0.9)$ with normal
$(-0.5, 0.866)$, with $v_{\max} = 1.5$ and $\vel^{\mathrm{pref}} = (1.2, 0.3)$.
(a) Run \cref{alg:ch13-lp} by hand in the order $H_1, H_2, H_3$ and in the
order $H_3, H_2, H_1$; report the calls to \OrcaLPLine and the result.
(b) Prove the claim of the walkthrough: if the optimum for the first $i-1$
half-planes violates $H_i$, the optimum for the first $i$ half-planes lies
on the boundary line of $H_i$. (c) Now let $\vel^{\mathrm{pref}} = (1.6, 0.4)$
and show where the disc enters the computation.
\end{exercise}

\begin{exercise}[\difficulty{2}]
\label{exr:ch13-infeasible}
Construct three half-planes whose intersection with the disc
$\norm{\vel} \le 1.5$ is empty, for instance three lines forming a triangle
that lies outside the disc with the permitted sides facing away from it.
(a) Solve \cref{eq:ch13-dense} for your instance geometrically: the solution
equalises the violations of the active half-planes. (b) Walk through
\cref{alg:ch13-dense} and identify the bisector half-planes it constructs.
(c) Explain why the result does not depend on $\vel^{\mathrm{pref}}$ and
whether that is desirable for a drone.
\end{exercise}

\begin{exercise}[\difficulty{2}]
\label{exr:ch13-corridor}
Two agents of radius $0.5$ face each other in a corridor of width $1.5$
along the $x$-axis, walls at $y = \pm 0.75$ represented by discs of radius
$0.5$ spaced $0.2$ apart with their centres at $y = \pm 1.25$. (a) Write down
the half-planes that the two nearest wall discs induce on an agent flying
$(1, 0)$ at $y = 0$ with $\ttc_{\mathrm{obst}} = 2$ and show that together
they bound $\abs{v_y}$ by a fixed multiple of $v_x$ that you compute.
(b) Show that in the head-on arc case each agent may close the gap
$d = \norm{\pos_{\mathrm{rel}}} - R$ at no more than $d/(2\ttc)$, derive
$d_{k+1} = d_k(1 - \dt/\ttc)$ for the gap after one step of
\cref{sec:ch13-step}, and conclude that the configuration converges to a
stalled state instead of reaching it in finitely many steps. (c) Propose two remedies, one
inside the local layer (a share or priority rule) and one in the
architecture of \cref{ch:ch24}, and discuss what each guarantees. Simulate
the instance with \code{ch13\_orca.py} if you want to see the stall.
\end{exercise}

\begin{exercise}[\difficulty{3} Coding]
\label{exr:ch13-coding}
This is the coding exercise of Week~6. Build a 2D simulation with two
groups of four agents whose straight paths cross at right angles
(\code{crossing\_scenario} in \code{ch13\_orca.py} gives the instance) and
compare no avoidance, the velocity obstacle rule and \orca. Report for each
the number of colliding pairs, the smallest separation, the last arrival
time, the mean path length and the number of velocity reversals, as in
\cref{tab:ch13-circle}; plot the traces and the smallest separation over time
as in \cref{fig:ch13-circle}. Then vary $\ttc \in \set{2, 5, 10}$ and the
group spacing, and describe when the \orca program becomes infeasible and
what the dense fallback does then. Finally add an intruder that flies
straight through the crossing without avoiding, treat it with full
responsibility, and check that no swarm agent collides with it.
\end{exercise}

\begin{exercise}[\difficulty{3} Coding]
\label{exr:ch13-3d}
Use \code{orca\_half\_space\_3d} to build the half-spaces of six drones that
start on a sphere of radius $10$ and fly to their antipodal points, solve
the three-dimensional program by sampling velocities in the ball
$\norm{\vel} \le v_{\max}$ (or by the recursive incremental algorithm), and
verify that the smallest centre distance never drops below $r_A + r_B$.
\end{exercise}
